% !TEX root = MosicaFE_Guidebook.tex
% ===========================================================================
%  MosicaFE_Guidebook.tex
%
%  MosicaFE 使用说明（按 guidebook.pdf 的模板撰写）
%
%  编译方式（推荐 xelatex，因为正文含中文）：
%      xelatex MosicaFE_Guidebook.tex
%      xelatex MosicaFE_Guidebook.tex      % 第二遍生成目录与交叉引用
%
%  或者：latexmk -xelatex MosicaFE_Guidebook.tex
% ===========================================================================
\documentclass[UTF8,11pt,oneside]{ctexbook}

\usepackage{amsmath}
\usepackage{amssymb}
\usepackage{amsthm}
\usepackage{booktabs}
\usepackage{longtable}
\usepackage{array}
\usepackage{graphicx}
\usepackage{listings}
\usepackage{upquote}
\usepackage{xcolor}
\usepackage{tcolorbox}
\usepackage{algorithm}
\usepackage{algorithmic}
\usepackage{enumitem}
\usepackage{caption}
\usepackage{geometry}
\usepackage{fancyhdr}
\usepackage[hidelinks]{hyperref}

\geometry{a4paper,left=2.6cm,right=2.6cm,top=2.8cm,bottom=2.8cm}

% ---------------------------------------------------------------------------
% 代码排版
% ---------------------------------------------------------------------------
\definecolor{codebg}{RGB}{248,248,248}
\definecolor{codegreen}{RGB}{0,110,60}
\definecolor{codegray}{RGB}{120,120,120}
\definecolor{codeblue}{RGB}{20,60,160}
\definecolor{codered}{RGB}{170,30,30}

\lstset{
  basicstyle=\ttfamily\small,
  keywordstyle=\color{codeblue}\bfseries,
  commentstyle=\color{codegreen}\itshape,
  stringstyle=\color{codered},
  numberstyle=\tiny\color{codegray},
  backgroundcolor=\color{codebg},
  frame=single,
  rulecolor=\color{codegray!40},
  framexleftmargin=4pt,
  numbers=left,
  numbersep=8pt,
  showstringspaces=false,
  breaklines=true,
  breakatwhitespace=false,
  tabsize=4,
  columns=fullflexible,
  keepspaces=true,
  captionpos=b,
  xleftmargin=18pt,
  framexrightmargin=0pt,
  upquote=true
}

\lstdefinelanguage{MosicaFE}{
  language=Python,
  morekeywords={as,with,None,True,False,lambda},
}

% 控制台输出样式
\lstdefinestyle{shell}{
  language=,
  numbers=none,
  backgroundcolor=\color{black!4},
  basicstyle=\ttfamily\footnotesize
}
\lstdefinestyle{textout}{
  language=,
  numbers=none,
  backgroundcolor=\color{black!4},
  basicstyle=\ttfamily\footnotesize
}

% ---------------------------------------------------------------------------
% 定理类环境
% ---------------------------------------------------------------------------
\theoremstyle{definition}
\newtheorem{definition}{定义}[chapter]
\newtheorem{example}{例}[chapter]
\newtheorem{remark}{注}[chapter]
\newtheorem{property}{性质}[chapter]
\theoremstyle{plain}
\newtheorem{theorem}{定理}[chapter]
\newtheorem{lemma}{引理}[chapter]

\newenvironment{important}{%
  \begin{tcolorbox}[colback=blue!5,colframe=blue!45,title=要点,fonttitle=\bfseries]%
}{%
  \end{tcolorbox}%
}

\newenvironment{caution}{%
  \begin{tcolorbox}[colback=orange!5,colframe=orange!55,title=注意,fonttitle=\bfseries]%
}{%
  \end{tcolorbox}%
}

\newcommand{\mosica}{\textsf{MosicaFE}}
\newcommand{\code}[1]{\texttt{#1}}
\newcolumntype{L}[1]{>{\raggedright\arraybackslash}p{#1}}
\newcommand{\module}[1]{\texttt{MosicaFE.#1}}
\newcommand{\dif}{\,\mathrm{d}}

\captionsetup{font=small,labelfont=bf}
\renewcommand{\arraystretch}{1.12}

\title{\Huge \textbf{MosicaFE}\\[0.4em]
       \Large 统一调用 FEM、VEM 与 RT0 求解偏微分方程的 Python 库\\[0.6em]
       \large 使用手册与开发指南}
\author{\normalsize MosicaFE contributors}
\date{\normalsize 2026 年 9 月 20 日}

% ===========================================================================
\begin{document}

\frontmatter
\maketitle

\begin{center}
\begin{minipage}{0.86\textwidth}
\small
\textbf{摘要.} 本手册介绍 \mosica{}——一个把有限元（FEM）、虚拟元（VEM）与
最低阶 Raviart--Thomas 混合元（RT0$\times$P0）统一到同一套接口下的 Python 库。
用户只需在主文件里写几行代码：选择网格、选择"元"、选择线性求解器、选择可视化工具，
即可离散并求解 Poisson 型偏微分方程。库的分层设计（\module{core}、\module{spaces}、
\module{physics}、\module{solvers}、\module{postprocess}）使得\textbf{新增一种单元}或
\textbf{新增一门方程}都是局部改动。手册给出完整的调用示例、数学背景、精度与收敛性验证，
以及扩展框架的模板。
\end{minipage}
\end{center}

\vspace{1.0em}

\tableofcontents

\mainmatter

% ===========================================================================
\chapter{引言}
\label{chap:intro}

\section{概述}

\mosica{} 是面向教学与科研的偏微分方程求解库。它把原本互相独立的三套程序
（经典有限元、虚拟元、RT0 混合有限元）整合成\textbf{一个包、一套接口}，
并使用同一套网格、求解器与后处理工具。

库的核心设计目标是：

\begin{enumerate}[label=(\arabic*),itemsep=0.25em]
  \item \textbf{调用极简}：主文件通常不超过 10 行；
  \item \textbf{模块解耦}：网格、单元、方程、线性代数、后处理各自独立；
  \item \textbf{易于扩展}：新增单元或新增 PDE 都只需实现少量接口；
  \item \textbf{可验证}：内置收敛性研究与误差范数，收敛阶可一键复现；
  \item \textbf{自包含}：只依赖 \code{numpy}（\code{scipy} 与 \code{matplotlib} 可选）。
\end{enumerate}

\section{设计哲学}

\mosica{} 的设计对应有限元方法的数学结构：

\begin{equation}
\underbrace{\text{PDE}}_{\text{physics}}
\;\xrightarrow{\ \text{弱形式}\ }\;
\underbrace{\text{变分问题}}_{\text{spaces}}
\;\xrightarrow{\ \text{离散化}\ }\;
\underbrace{Au=b}_{\text{assembly}}
\;\xrightarrow{\ \text{求解}\ }\;
\underbrace{u_h}_{\text{postprocess}}
\label{eq:philosophy}
\end{equation}

每一步由独立的模块负责，模块之间只通过明确的接口通信：

\begin{itemize}[itemsep=0.2em]
  \item \module{core} 提供网格与求积规则；
  \item \module{spaces} 把"选哪种单元"翻译成\textbf{矩阵级接口}：
        刚度矩阵、质量矩阵、载荷向量（详见第 \ref{chap:spaces} 章）；
  \item \module{physics} 只做这些矩阵的线性组合，因此描述一门新 PDE 通常只要十几行；
  \item \module{solvers} 负责线性代数；
  \item \module{postprocess} 负责误差、收敛性与可视化。
\end{itemize}

\begin{important}
"元"与"方程"的正交分解是 \mosica{} 与三套原始程序最大的区别。
原程序把"三角形网格 + 线性元 + Poisson + \code{spsolve} + 画图"
写死在同一个文件里；现在这几件事互不相干：
换一种元不必改方程，换一门方程不必改单元。
\end{important}

\section{快速开始}

下面是一个完整的 2D Poisson 求解程序：

\begin{lstlisting}[language=MosicaFE,caption={示例 1.1：两行求解 Poisson 方程},label={lst:quickstart}]
import numpy as np
from MosicaFE import Mesh, solve_poisson

source = lambda x: 2*np.pi**2 * np.sin(np.pi*x[0]) * np.sin(np.pi*x[1])

u = solve_poisson(Mesh.rectangle(nx=32, ny=32), element="P2", f=source)
print("deg of freedom:", u.size)
\end{lstlisting}

等价的"分步"写法（也是本手册后续使用的标准形式）：

\begin{lstlisting}[language=MosicaFE,caption={示例 1.2：选择网格 / 元 / 求解器 / 可视化},label={lst:quicksplit}]
import numpy as np
from MosicaFE import Mesh, FESpace, PoissonProblem
from MosicaFE.postprocess import plot_solution

def exact(x):
    return np.sin(np.pi*x[0]) * np.sin(np.pi*x[1])

def source(x):
    return 2*np.pi**2 * exact(x)

mesh = Mesh.rectangle(nx=32, ny=32)         # (1) 选网格
V = FESpace.lagrange(mesh, degree=2)        # (2) 选元
problem = PoissonProblem(V, f=source, g=lambda x: 0.0)
u = problem.solve(method="direct")          # (3) 选求解器
plot_solution(u, mesh, V, filename="p2.png")  # (4) 选可视化

print(problem.compute_errors(u, exact))
\end{lstlisting}

把第 \ref{lst:quicksplit} 行中的 \code{FESpace.lagrange(mesh, degree=2)}
换成 \code{FESpace.vem(mesh)}，得到的就不再是有限元，
而是\textbf{同一个方程}在\textbf{同一种网格}上的虚拟元解；程序其它部分一字不改。

\begin{important}
这就是本库存在的意义：让"换方法"的代价降到一行代码。
\end{important}

\section{包结构}
\label{sec:package-structure}

\begin{table}[htbp]
\centering
\caption{\mosica{} 的模块划分}
\label{tab:modules}
\begin{tabular}{@{}lp{10.4cm}@{}}
\toprule
模块 & 职责 \\
\midrule
\module{core} & 网格数据结构与生成（区间 / 矩形 / 长方体 / 多边形 / Voronoi）、
\textbf{面拓扑与定向}（\code{facets.py}，$H(\mathrm{div})$ 协调空间的基础）、
求积规则、几何工具（面积、体积、雅可比、梯度变换） \\
\module{spaces} & 离散空间：\code{LagrangeSpace}（P1/P2/Q1/Q2）、
\code{VEMSpace}（$k=1$ 虚拟元）、\code{RT0Space}（RT0$\times$P0），
以及工厂 \code{FESpace} \\
\module{physics} & PDE 弱形式与组装：\code{PoissonProblem}、
\code{MixedPoissonProblem}、\code{ReactionDiffusionProblem}，含 PDE 注册表 \\
\module{solvers} & 线性求解器：直接法、CG、GMRES、BiCGSTAB、MINRES 与自动选择 \\
\module{postprocess} & 误差范数（$L^2$ / $H^1$ / $L^\infty$）、
收敛性研究、可视化、混合格式（压力 + 通量）的误差范数 \\
\bottomrule
\end{tabular}
\end{table}

\begin{important}
FEM、VEM、RT0 三类元\textbf{全部是库自身的实现}：RT0$\times$P0 的显式基函数、
精确质量矩阵、面拓扑与鞍点组装分别位于 \code{spaces/rt0.py}、
\code{core/facets.py} 与 \code{physics/mixed\_poisson.py}，
\mosica{} 不依赖、也不需要接入任何外部有限元后端。
\end{important}

\section{本手册的组织}

\begin{itemize}[itemsep=0.2em]
  \item 第 \ref{chap:install} 章：安装与环境；
  \item 第 \ref{chap:math} 章：三种方法的数学背景；
  \item 第 \ref{chap:core}--\ref{chap:spaces} 章：网格、求积与离散空间；
  \item 第 \ref{chap:physics}--\ref{chap:solvers} 章：方程组装与线性求解；
  \item 第 \ref{chap:error}--\ref{chap:visual} 章：误差分析与可视化；
  \item 第 \ref{chap:examples} 章：完整示例；
  \item 第 \ref{chap:advanced}--\ref{chap:extend} 章：进阶话题与框架扩展；
  \item 第 \ref{chap:faq} 章：常见问题；
  \item 附录 \ref{app:api}--\ref{app:bib}：API 速查、数学符号、参考文献。
\end{itemize}

% ===========================================================================
\chapter{安装与环境}
\label{chap:install}

\section{系统要求}

\subsection{Python 版本}

要求 Python $\ge$ 3.9。开发与验证使用的环境为
Python 3.12.10（conda 环境 \code{fealpy}），\code{numpy} 2.2.4、
\code{scipy} 1.15.2、\code{matplotlib} 3.10.1。

\subsection{依赖}

\begin{table}[htbp]
\centering
\caption{依赖关系}
\label{tab:deps}
\begin{tabular}{@{}lll@{}}
\toprule
包 & 是否必需 & 用途 \\
\midrule
\code{numpy} & 必需 & 数组、稠密线性代数 \\
\code{scipy} & 可选 & 稀疏矩阵、\code{spsolve}、Krylov 求解器、Voronoi 网格 \\
\code{matplotlib} & 可选 & 可视化 \\
\bottomrule
\end{tabular}
\end{table}

缺少 \code{scipy} 时，\mosica{} 自动退化为纯 \code{numpy} 的稠密实现
（\module{\_compat} 统一处理），小规模问题照常可解；
缺少 \code{matplotlib} 时，所有绘图函数打印一条提示并返回 \code{None}，不会报错。

\section{安装}

\subsection{方式一：以包的形式安装（推荐）}

\begin{lstlisting}[style=shell]
cd D:\mosica\MosicaFE
python -m pip install -e .
\end{lstlisting}

\subsection{方式二：直接使用源码}

把 \mosica{} 的\textbf{父目录}加入模块搜索路径即可，无需安装：

\begin{lstlisting}[style=shell]
# PowerShell
cd D:\mosica
$env:PYTHONPATH = "D:\mosica"
& "D:\anaconda\envs\fealpy\python.exe" MosicaFE\examples\poisson_demo.py
\end{lstlisting}

\section{验证安装}

\mosica{} 自带两个不依赖第三方测试框架的脚本：

\begin{lstlisting}[style=shell]
& "D:\anaconda\envs\fealpy\python.exe" MosicaFE\tests\run_tests.py
& "D:\anaconda\envs\fealpy\python.exe" MosicaFE\tests\run_checks.py
\end{lstlisting}

前者跑 54 个单元/集成测试，后者打印各方法的实测收敛阶。预期输出片段：

\begin{lstlisting}[style=textout]
=== P2 ===
  level        h                  L2    rate          H1    rate        Linf    rate
      0  0.35355    3.4643e-02       -  4.4391e-01       -  1.4645e-01       -
      1  0.17678    4.4076e-03    2.97  1.1632e-01    1.93  3.8060e-02    1.94
      2  0.08839    5.5040e-04    3.00  2.9416e-02    1.98  9.6074e-03    1.99
      3  0.04419    6.8699e-05    3.00  7.3752e-03    2.00  2.4076e-03    2.00
  rate[L2] (last 3) = 3.002
  rate[H1] (last 3) = 1.990
\end{lstlisting}

如果安装了 \code{pytest}，也可以直接：

\begin{lstlisting}[style=shell]
python -m pytest MosicaFE/tests -q
\end{lstlisting}

\section{编译本手册}

本文档使用 \code{ctexbook} 文档类，正文含中文，需要用 XeLaTeX{} 编译两遍：

\begin{lstlisting}[style=shell]
xelatex MosicaFE_Guidebook.tex
xelatex MosicaFE_Guidebook.tex
\end{lstlisting}

在 TeXShop 中把引擎选为 \code{XeLaTeX} 并点击 Typeset 即可。

% ===========================================================================
\chapter{数学背景}
\label{chap:math}

\section{模型问题}

本库（以及原来的三套程序）求解的核心模型问题是 Poisson 方程

\begin{equation}
\begin{cases}
-\nabla\cdot(\kappa\nabla u) = f, & \text{in } \Omega,\\[2pt]
u = g, & \text{on } \partial\Omega,
\end{cases}
\label{eq:poisson}
\end{equation}

其中 $\Omega\subset\mathbb{R}^d$（$d=1,2,3$）是有界区域，
$\kappa>0$ 是扩散系数。当 $g=0$ 且 $\kappa=1$ 时即经典的
$-\Delta u=f$。

\section{弱形式}

\begin{definition}[弱形式]
记 $V=\{v\in H^1(\Omega): v|_{\partial\Omega}=0\}$，
$V_g=\{v\in H^1(\Omega): v|_{\partial\Omega}=g\}$。
求 $u\in V_g$ 使得
\begin{equation}
a(u,v)=\ell(v),\qquad \forall v\in V,
\label{eq:weak}
\end{equation}
其中
\begin{equation}
a(u,v)=\int_\Omega \kappa\nabla u\cdot\nabla v\dif x,\qquad
\ell(v)=\int_\Omega f v\dif x .
\label{eq:bilinear}
\end{equation}
\end{definition}

$\ell$ 也可以带 Neumann 数据项 $\int_{\Gamma_N} h v\dif s$，
\mosica{} 的物理层接口为此预留了扩展点。

\section{有限元离散（FEM）}

设 $V_h\subset V$ 是有限维子空间，基函数为 $\{\varphi_i\}_{i=1}^n$。
离散问题为：求 $u_h=\sum_j u_j\varphi_j$ 使得

\begin{equation}
\sum_{j=1}^n u_j\,a(\varphi_j,\varphi_i)=\ell(\varphi_i),\qquad i=1,\dots,n,
\label{eq:fem-system}
\end{equation}

即线性方程组
\begin{equation}
A u = b,\qquad
A_{ij}=\int_\Omega\kappa\nabla\varphi_j\cdot\nabla\varphi_i\dif x,\qquad
b_i=\int_\Omega f\varphi_i\dif x .
\label{eq:Au}
\end{equation}

在参考单元 $\hat K$ 上计算时用到梯度变换

\begin{equation}
\nabla_x\varphi_i = J_K^{-\top}\nabla_{\hat x}\hat\varphi_i,
\qquad
\int_K g\dif x=\int_{\hat K}g\circ F_K\,|\det J_K|\dif\hat x .
\label{eq:transform}
\end{equation}

\begin{table}[htbp]
\centering
\caption{\mosica{} 中 Lagrange 元的信息}
\label{tab:lagrange}
\begin{tabular}{@{}llcc@{}}
\toprule
单元 & 名字 & 每个单元的局部自由度 & 自由度类型 \\
\midrule
三角形 & P1 & 3 & 顶点 \\
三角形 & P2 & 6 & 顶点 + 棱中点 \\
四边形 & Q1 & 4 & 顶点 \\
四边形 & Q2 & 9 & 顶点 + 棱中点 + 单元中心 \\
四面体 & P1 & 4 & 顶点 \\
四面体 & P2 & 10 & 顶点 + 棱中点 \\
六面体 & Q1 & 8 & 顶点 \\
六面体 & Q2 & 27 & 顶点 + 棱 + 面 + 体 \\
\bottomrule
\end{tabular}
\end{table}

\section{虚拟元离散（VEM）}
\label{sec:vem-math}

虚拟元方法的关键思想是：\textbf{不需要写出基函数的显式表达式}。
一次（$k=1$）虚拟元空间定义为

\begin{equation}
V_h(K)=\Bigl\{v\in H^1(K):\ \Delta v\in\mathcal{P}_1(K),\
v|_e\in\mathcal{P}_1(e)\ \forall e\subset\partial K\Bigr\},
\label{eq:vem-space}
\end{equation}

其自由度就是每个顶点上的函数值。单元上的双线性型被拆成两部分：

\begin{equation}
a_h(u_h,v_h)=\underbrace{a\bigl(\Pi^{\nabla}u_h,\Pi^{\nabla}v_h\bigr)}_{A_{\rm proj}}
+\underbrace{S_K\bigl((I-\Pi)u_h,(I-\Pi)v_h\bigr)}_{S_{\rm stab}} .
\label{eq:vem-split}
\end{equation}

\subsection{投影算子}

取缩放单项式基 $m_0=1$，$m_{k+1}=x_k-x_{K,k}$（$x_K$ 为单元形心）。
投影 $\Pi^{\nabla}:V_h(K)\to\mathcal{P}_1(K)$ 满足

\begin{equation}
\int_K\nabla\Pi^{\nabla}v\cdot\nabla q\dif x
=\int_K\nabla v\cdot\nabla q\dif x
=-\int_K v\,\Delta q\dif x+\int_{\partial K}v\,(\nabla q\cdot n)\dif s
=\int_{\partial K}v\,(\nabla q\cdot n)\dif s ,
\label{eq:vem-proj}
\end{equation}

最后一步用到 $\Delta q=0$（$q$ 是一次多项式）。
由于 $v|_e$ 在每条边上是一次函数，
$\int_e v\,(\nabla q\cdot n)\dif s$ 用\textbf{梯形法则即精确}：

\begin{equation}
\int_e v\,(\nabla q\cdot n)\dif s
=\frac{|e|}{2}\Bigl(v(a)(\nabla q\cdot n)+v(b)(\nabla q\cdot n)\Bigr).
\label{eq:vem-edge}
\end{equation}

写成矩阵形式：设 $G=[\nabla m_0,\dots,\nabla m_d]\in\mathbb{R}^{d\times(d+1)}$，
$M=|K|\,G^{\top}G$，$R_{\alpha i}=\int_{\partial K}\varphi_i(\nabla m_\alpha\cdot n)\dif s$，
则

\begin{equation}
\Pi=\text{pinv}(M)\,R\in\mathbb{R}^{(d+1)\times n}.
\label{eq:vem-Pi}
\end{equation}

\begin{caution}
常函数模式的规范化至关重要。若取 $\Pi_{0i}=1/n$（顶点值的算术平均），
则只有当单元"顶点算术平均 $=$ 形心"时（三角形、正方形成立，
一般多边形不成立）投影才精确再生一次多项式。
正确的规范化由
$\ p(x_K)=\overline{d}-\nabla p\cdot(\overline{X}-x_K)$ 给出：

\begin{equation}
\Pi_{0\cdot}=\frac{1}{n}\mathbf{1}^{\top}
-(\overline{X}-x_K)\,\Pi_{1:d,\cdot},
\qquad \overline{X}=\frac{1}{n}\sum_i X_i .
\label{eq:vem-const}
\end{equation}

若忽略这一点，稳定化项会惩罚线性函数，相容性被破坏，
在不规则多边形网格上 $L^2$ 收敛阶会从 2 掉到 0.5 以下。
细节见第 \ref{sec:pitfall-vem} 节。
\end{caution}

\subsection{稳定化项}

\mosica{} 默认采用\textbf{对称半正定}的稳定化

\begin{equation}
S_{\rm stab}=\sigma\,(I-\Pi_{\rm dof})^{\top}(I-\Pi_{\rm dof}),\qquad
\Pi_{\rm dof}=D\,\Pi,\quad
D_{i\cdot}=[1,\ X_i-x_K],
\label{eq:vem-stab}
\end{equation}

其中缩放因子取

\begin{equation}
\sigma=\frac{\operatorname{tr}(A_{\rm proj})}{\operatorname{tr}\bigl((I-\Pi_{\rm dof})^{\top}(I-\Pi_{\rm dof})\bigr)}
\label{eq:vem-sigma}
\end{equation}

使两部分量级相当（也自动对齐量纲）。
传 \code{stabilization="legacy"} 可以复现早期脚本中
$(I-\Pi_{\rm dof})/n$ 的形式（非对称，不推荐）；
\code{stabilization="none"} 只用于诊断，因为纯投影项是奇异的。

\subsection{载荷向量与质量矩阵}

虚拟基函数没有显式表达式，因此 $\int_K f\varphi_i$ 需要用\textbf{扇形三角剖分
（3D 为四面体剖分）上的分片线性插值}近似。把单元从形心剖成
$T_1,\dots,T_n$，每个 $T_i$ 的三个角分别为形心 $c$ 与 $V_i,V_{i+1}$，
其重心坐标为 $L_0,L_1,L_2$，则

\begin{equation}
\int_K f\varphi_j\dif x
\approx c_0[j]\sum_i\int_{T_i} f L_0\dif x
+\int_{T_j}f L_1\dif x
+\int_{T_{j-1}}f L_2\dif x,
\label{eq:vem-load}
\end{equation}

其中 $c_0[j]$ 是基函数 $\varphi_j$ 在形心处的插值
（取线性最小二乘重构的值，既保证线性再生又保证单位分解）。
对三角形网格该公式是精确的；对一般多边形它保留了正确的几何权重。

\begin{caution}
早期脚本对所有顶点统一用 $|K|/n$ 作为 $\int_K\varphi_i$。
在四边形和三角形上这个近似恰好足够好，但在不规则多边形上
会严重高估"短边"上的顶点，导致收敛停滞。这是我们把载荷与质量矩阵
改为公式 \eqref{eq:vem-load} 的直接原因。
\end{caution}

\section{RT0$\times$P0 混合离散}

混合形式把 Poisson 方程写成一阶系统

\begin{equation}
q=-\nabla p,\qquad \nabla\cdot q=f\ \text{ in }\Omega,\qquad
p=0\ \text{ on }\partial\Omega .
\label{eq:mixed}
\end{equation}

弱形式：求 $(q,p)\in RT0\times P0$ 使得

\begin{equation}
\int_\Omega \nu\cdot q\dif x-\int_\Omega p\,\nabla\cdot\nu\dif x=0,
\qquad
\int_\Omega r\,\nabla\cdot q\dif x=\int_\Omega f r\dif x
\label{eq:mixed-weak}
\end{equation}

对任意 $(\nu,r)$ 成立。离散后得到鞍点系统

\begin{equation}
\begin{pmatrix} M & B^{\top}\\ B & 0\end{pmatrix}
\begin{pmatrix} \mathbf{Q}\\ \mathbf{P}\end{pmatrix}
=\begin{pmatrix} 0\\ \mathbf{f}\end{pmatrix},
\label{eq:saddle}
\end{equation}

它是\textbf{对称不定}的，因此不能使用共轭梯度法。

\begin{table}[htbp]
\centering
\caption{RT0$\times$P0 的自由度和边界条件}
\label{tab:rt0}
\begin{tabular}{@{}lll@{}}
\toprule
量 & 自由度位置 & 个数 \\
\midrule
通量 $q_h$ & 每个单元面（2D 边 / 3D 三角面），相邻单元共享 & $\#$facets \\
压力 $p_h$ & 每个单元一个常数 & $\#$cells \\
\midrule
Dirichlet 条件 $p=0$ & \multicolumn{2}{l}{混合格式的\textbf{自然}边界条件，无需消去} \\
\bottomrule
\end{tabular}
\end{table}

\section{误差估计}

\begin{theorem}[Cea 引理]
设 $u$ 是精确解、$u_h$ 是有限元解，则
$\|u-u_h\|_V\le C\inf_{v_h\in V_h}\|u-v_h\|_V$。
\end{theorem}

\begin{theorem}[先验误差估计]
若 $u\in H^{k+1}(\Omega)$、采用 $k$ 次多项式，则
\begin{equation}
\|u-u_h\|_{L^2}\le C h^{k+1}|u|_{H^{k+1}},\qquad
|u-u_h|_{H^1}\le C h^{k}|u|_{H^{k+1}} .
\label{eq:apriori}
\end{equation}
\end{theorem}

\begin{example}[本库实测的收敛阶]
\begin{itemize}[itemsep=0.15em]
  \item P1（三角形）：$L^2$ 为 $O(h^2)$，$H^1$ 为 $O(h)$；
  \item P2 / Q2：$L^2$ 为 $O(h^3)$，$H^1$ 为 $O(h^2)$；
  \item VEM（$k=1$，任意多边形）：$L^2$ 为 $O(h^2)$，$H^1$ 为 $O(h)$；
  \item RT0$\times$P0：压力 $L^2$ 为 $O(h)$，通量 $L^2$ 为 $O(h)$。
\end{itemize}
第 \ref{chap:error} 章给出实测数据。
\end{example}

% ===========================================================================
\chapter{核心模块：网格与求积}
\label{chap:core}

\section{网格数据结构}

\subsection{网格的组成部分}

\code{Mesh} 对象保存以下信息：

\begin{table}[htbp]
\centering
\caption{\code{Mesh} 的主要属性}
\label{tab:mesh-attrs}
\begin{tabular}{@{}L{2.8cm}L{1.9cm}L{9.9cm}@{}}
\toprule
属性 & 类型 & 含义 \\
\midrule
\code{nodes} & \code{(n\_nodes, dim)} 数组 & 顶点坐标 \\
\code{elements} & \code{list[ndarray]} & 每个单元的顶点编号，\textbf{变长} \\
\code{element\_type} & \code{str} & \code{interval}/\code{triangle}/\code{quad}/\code{tet}/\code{hex}/\code{polygon}/\code{polyhedron} \\
\code{boundary\_nodes} & \code{ndarray} & 边界顶点编号（由 facet 拓扑自动判定） \\
\code{dim} & \code{int} & 空间维数 \\
\bottomrule
\end{tabular}
\end{table}

\begin{important}
\code{elements} 故意设计成\textbf{变长列表}：三角形、四边形、五边形、
Voronoi 多边形、四面体、六面体因此可以共用同一套组装循环。
如果所有单元顶点数相同，可以通过 \code{mesh.uniform\_vertices}
取出 \code{(n\_cells, k)} 的数组视图。
\end{important}

\subsection{facet 拓扑}

余维 1 的单元面（facet）是边界判定的基础：
1D 的 facet 是端点、2D 是边、3D 是三角形或四边形面。
\code{mesh.facets} 给出全部 facet，
\code{mesh.boundary\_facets} 给出只被一个单元拥有的 facet，
它们的顶点集合就是 Dirichlet 边界所在的位置。

\begin{property}
对二维任意多边形网格，facet 就是多边形的边；
对六面体网格按固定的 6 个面表展开；
对一般三维多面体，可以通过 \code{Mesh(..., faces=...)} 显式提供面表。
\end{property}

\begin{caution}
上面这套 \code{mesh.facets} 是"\textbf{边界判定}用的顶点集合"，
它足以支撑 Lagrange 与虚拟元（自由度在顶点上）。
但对 $H(\mathrm{div})$ 协调的混合元（RT0）来说，还额外需要
"\textbf{哪些单元共享同一个面}"以及"这个面的全局法向指向谁"，
因此本库另有一层面拓扑：\code{core/facets.py}。下一小节给出它的用法。
\end{caution}

\subsection{面拓扑与定向：\code{core/facets.py}}
\label{sec:facets}

RT0 的自由度是\textbf{面法向通量}
$\Phi_f=\int_f q\cdot n_f\dif s$，相邻两个单元必须共享同一个未知量。
这要求显式地给出面的全局编号与定向。本库把它做成 \code{core} 层的一等公民
\code{FacetTopology}（\code{mesh.facet\_topology} 惰性构建并缓存）：

\begin{table}[htbp]
\centering
\caption{\code{FacetTopology} 的主要属性（\code{t = mesh.facet\_topology}）}
\label{tab:facettopo}
\begin{tabular}{@{}L{4.4cm}L{10.0cm}@{}}
\toprule
属性 / 方法 & 含义 \\
\midrule
\code{t.facets} & \code{(n\_facets, d)}：每个面的\textbf{升序}顶点表（全局唯一键） \\
\code{t.cell\_facets} & \code{(n\_cells, d+1)}：\code{[c,i]} 是单元 $c$ 第 $i$ 个局部面的全局编号 \\
\code{t.signs} & $\pm1$：全局法向是否指向该单元外部 \\
\code{t.normals} & 未归一化的全局法向（模长 = 面测度） \\
\code{t.counts} & 相邻单元个数：内部面 2、边界面 1 \\
\code{t.boundary\_facets()} & 边界面编号数组（\code{t.n\_boundary} 为个数） \\
\bottomrule
\end{tabular}
\end{table}

\begin{lstlisting}[language=MosicaFE,caption={示例 4.2：面拓扑与定向}]
from MosicaFE import Mesh

mesh = Mesh.box(nx=2, ny=2, nz=2, element_type="tet")
t = mesh.facet_topology

print(t.summary())          # cells=48, facets=120, interior=72, boundary=48
print(t.cell_facets[0])     # 单元 0 的 4 个局部面的全局编号
print(t.signs[0])           # 外法向符号（+1 / -1）

# 内部面的两侧符号一定相反 —— 这就是 H(div) 协调性的离散体现
c0, c1 = t.facet_cells[0]
print(t.signs[c0][list(t.cell_facets[c0]).index(0)] *
      t.signs[c1][list(t.cell_facets[c1]).index(0)])   # -1
\end{lstlisting}

\begin{important}
\code{build\_facet\_topology} 会拒绝两类"会静默算错"的输入：
非单纯形网格（四边形 / 多边形 / 六面体），
以及存在被 3 个以上单元共用的面（不协调剖分）。
这是 RT0 只接受三角形 / 四面体网格的技术原因。
\end{important}

\section{网格生成}
\label{sec:mesh-generation}

所有网格都由 \code{Mesh} 的类方法生成，且统一采用
\textbf{逆时针（2D）/ 正体积（3D）}取向。

\begin{lstlisting}[language=MosicaFE,caption={示例 4.1：生成各类网格},label={lst:meshgen}]
from MosicaFE import Mesh

m1 = Mesh.interval(a=0.0, b=1.0, nx=20)                # 1D interval

m2 = Mesh.rectangle(nx=8, ny=8)                        # triangles
m3 = Mesh.rectangle(nx=8, ny=8, element_type="quad")   # quadrilaterals

m4 = Mesh.box(nx=4, ny=4, nz=4)                        # tets (Kuhn split)
m5 = Mesh.box(nx=4, ny=4, nz=4, element_type="hex")    # hexahedra
m6 = Mesh.box(nx=3, ny=3, nz=3, pattern="center")      # 12 tets per cube

m7 = Mesh.pentagon(nx=8, ny=8)                         # pentagon mesh
m8 = Mesh.voronoi_polygon(n_points=100, seed=42)       # Voronoi polygons

for m in (m1, m2, m3, m4, m5, m7, m8):
    print(m.summary())
\end{lstlisting}

\begin{table}[htbp]
\centering
\caption{网格工厂速查}
\label{tab:mesh-factories}
\begin{tabular}{@{}L{6.4cm}L{1.9cm}L{6.3cm}@{}}
\toprule
工厂 & 区域 & 单元 \\
\midrule
\code{Mesh.interval(a,b,nx)} & $[a,b]$ & 线段 \\
\code{Mesh.rectangle(...,"triangle")} & 矩形 & 三角形 \\
\code{Mesh.rectangle(...,"quad")} & 矩形 & 四边形 \\
\code{Mesh.box(...,"tet")} & 长方体 & 四面体 \\
\code{Mesh.box(...,"hex")} & 长方体 & 六面体 \\
\code{Mesh.pentagon(nx,ny)} & 单位正方形 & 五边形 \\
\code{Mesh.voronoi\_polygon(...)} & 单位正方形 & Voronoi 多边形 \\
\code{Mesh.from\_arrays(...)} & 自定义 & 自动推断 \\
\bottomrule
\end{tabular}
\end{table}

\begin{remark}[Kuhn 剖分]
\code{Mesh.box(...,element\_type="tet")} 默认把每个小立方体剖成 6 个
\textbf{共享主对角线}的四面体（Kuhn 剖分）。这样相邻立方体在公共面上的
三角剖分方向一致，网格是协调的。\code{pattern="center"} 则额外插入
立方体中心点，剖成 12 个四面体——与原始 VEM 三维脚本中的网格一致。
\end{remark}

\begin{remark}[Voronoi 网格]
Voronoi 网格生成时在区域外做 8 方向镜像反射点，
保证边界处的 Voronoi 单元闭合，然后裁剪回区域。
裁剪与排序之后会自动去掉相邻重合顶点，避免出现退化单元
（一个含重复顶点的多边形会让扇形剖分出现零面积三角形）。
\end{remark}

\section{网格属性与几何量}

\begin{lstlisting}[language=MosicaFE,caption={示例 4.2：读取网格信息}]
mesh = Mesh.rectangle(nx=8, ny=8)

print(mesh.dim, mesh.n_nodes, mesh.n_cells)     # 2 81 128
print(mesh.get_mesh_size())                     # max cell diameter
print(mesh.cell_measures())                     # array of areas

v = mesh.get_cell_vertices(0)                   # (k, dim) coordinates
h = mesh.get_cell_diameter(0)
c = mesh.get_cell_centroid(0)
print(mesh.boundary_nodes[:5])
\end{lstlisting}

\section{网格加密}

\begin{lstlisting}[language=MosicaFE,caption={示例 4.3：均匀加密}]
coarse = Mesh.rectangle(nx=4, ny=4)
fine = coarse.refine_uniform()
print(coarse.n_cells, fine.n_cells)     # 32 128
\end{lstlisting}

\code{refine\_uniform()} 只对结构化网格（\code{interval} /
\code{rectangle} / \code{box} / \code{pentagon}）有效：
它把各方向的网格数翻倍后重新生成，因此三角形网格的单元数变为 4 倍、
六面体网格变为 8 倍。对 Voronoi 这类非结构网格会抛出 \code{ValueError}。

\section{参考单元}

\begin{table}[htbp]
\centering
\caption{参考单元约定}
\label{tab:refelem}
\begin{tabular}{@{}ll@{}}
\toprule
单元类型 & 参考单元 \\
\midrule
\code{interval} & $[0,1]$ \\
\code{triangle} & $\{\xi\ge0,\eta\ge0,\xi+\eta\le1\}$ \\
\code{quad} & $[0,1]^2$，顶点顺序 $(0,0),(1,0),(1,1),(0,1)$ \\
\code{tet} & 标准四面体，顶点 $(0,0,0),(1,0,0),(0,1,0),(0,0,1)$ \\
\code{hex} & $[0,1]^3$，先 $z=0$ 面逆时针 4 点，再 $z=1$ 面对应 4 点 \\
\bottomrule
\end{tabular}
\end{table}

单纯形使用仿射映射 $x=v_0+J_K\hat x$；
四边形与六面体使用 $[0,1]^d$ 上的双线性 / 三线性映射，
其雅可比矩阵依赖参考点 $\hat x$，因此需要逐求积点计算。

\begin{lstlisting}[language=MosicaFE,caption={示例 4.4：几何工具}]
import numpy as np
from MosicaFE.core.utils import (
    compute_jacobian, transform_gradient, apply_affine_map,
    polygon_area, polyhedron_volume,
)

quad = np.array([[0,0],[2,0],[2,2],[0,2]], float)
J, detJ = compute_jacobian(quad, np.array([0.5, 0.5]))
print(J, detJ)                      # [[2,0],[0,2]] 4.0
print(polygon_area(quad))           # 4.0
print(apply_affine_map([[0.5,0.5]], quad))   # [[1, 1]]
\end{lstlisting}

\begin{caution}
\textbf{雅可比的约定。}\code{compute\_jacobian} 对两类单元返回\textbf{同一个约定}的雅可比：
列是物理坐标对参考坐标的偏导，

\begin{equation}
J_{ab}=\frac{\partial x_a}{\partial\xi_b},
\qquad
J=\big[\ \partial x/\partial\xi_1\ \ \cdots\ \ \partial x/\partial\xi_d\ \big],
\label{eq:jac-convention}
\end{equation}

只有这样，梯度变换 \code{transform\_gradient(dN, J)}（内部用 $J^{-\top}$）
与四边形 / 六面体的 Newton 逆映射（内部用 $J^{-1}$）才是对的。
\textbf{自己扩展代码时不要写成 $dN^{\top}V$}——那是转置雅可比。
对轴对齐的矩形 / 长方体，$J$ 是对角阵、转置不变，这个错误会被完全掩盖；
而网格一旦\textbf{旋转或畸变}（$J\ne J^{\top}$），梯度、刚度矩阵与插值
就会静默算错。本库在 2026-09-21 修掉过这一处（第 7 号 bug），
回归测试见 \code{tests/test\_tensor\_geometry.py}：它用有限差分核对 $J$，
并检验旋转 $0^\circ/15^\circ/30^\circ/45^\circ$ 时误差逐位相同。
\end{caution}

\section{求积规则}

\begin{lstlisting}[language=MosicaFE,caption={示例 4.5：使用求积规则}]
from MosicaFE.core.quadrature import QuadratureRule, cell_quadrature

quad = QuadratureRule.get_quadrature("triangle", order=3)
print(quad.n_points, quad.weights.sum())     # 4 0.5

f = lambda xi: xi[0]**2 + xi[1]**2
value = sum(quad.weights[q] * f(quad.points[q]) for q in range(quad.n_points))
print(value)                                  # 1/6

# 映射到物理单元（自动处理多边形 / 多面体的扇形剖分）
mesh = Mesh.voronoi_polygon(n_points=20, seed=1)
pts, w = cell_quadrature(mesh, 0, order=3)
print(w.sum(), mesh.get_cell_measure(0))
\end{lstlisting}

\begin{table}[htbp]
\centering
\caption{可用的求积规则（\code{order} 为需精确积分的多项式次数）}
\label{tab:quad}
\begin{tabular}{@{}lll@{}}
\toprule
单元 & 规则 & 说明 \\
\midrule
\code{interval} & Gauss--Legendre & $n=\lceil(order+1)/2\rceil$ 点，$\int_0^1$ 归一 \\
\code{triangle} & 重心 / 3 点 / 4 点 / 6 点 & 分别精确到 1/2/3/4 次 \\
\code{quad} & 张量 Gauss & 每个方向 $n$ 点 \\
\code{tet} & 1 点 / 4 点 / 5 点 & 精确到 1/2/3 次 \\
\code{hex} & 张量 Gauss & 每个方向 $n$ 点 \\
\bottomrule
\end{tabular}
\end{table}

\begin{caution}
张量积单元的求积阶需要按\textbf{每个方向}的次数来理解。
\code{Q1} 单元刚度矩阵的被积函数在每个方向上是 2 次，
所以每个方向至少需要 2 个 Gauss 点——\mosica{} 内部会自动把
阶数取成 \code{max(2, 2*p)}，用户一般不必手工指定。
早期版本在这里少取了一个点，导致三维六面体的收敛阶算出来是负数。
\end{caution}

% ===========================================================================
\chapter{离散空间}
\label{chap:spaces}

\section{空间抽象}

所有空间都继承 \code{BaseFESpace}，并向上提供\textbf{矩阵级接口}：

\begin{table}[htbp]
\centering
\caption{\code{BaseFESpace} 的关键方法}
\label{tab:space-api}
\begin{tabular}{@{}L{6.2cm}L{9.2cm}@{}}
\toprule
方法 & 含义 \\
\midrule
\code{stiffness\_matrix()} & $A_{ij}=\int\nabla\varphi_j\cdot\nabla\varphi_i$ \\
\code{mass\_matrix()} & $M_{ij}=\int\varphi_j\varphi_i$ \\
\code{load\_vector(f)} & $b_i=\int f\varphi_i$ \\
\code{boundary\_dofs()} & Dirichlet 自由度编号 \\
\code{interpolate(f)} & 把函数插值到自由度上 \\
\code{evaluate(cell,u,pts)} & 在单元内把自由度重构成函数并求值 \\
\code{evaluate\_gradient(cell,u,pts)} & 同上，返回梯度 \\
\code{nodal\_values(u)} & 把自由度向量还原到网格节点（可视化用） \\
\code{get\_cell\_dofs(i)} & 第 $i$ 个单元的全局自由度编号 \\
\code{get\_dof\_coordinates(k)} & 第 $k$ 个自由度的物理坐标 \\
\code{eval\_basis(xi)} / \code{eval\_basis\_grad(xi)} & 参考单元上的基函数（仅显式基函数空间） \\
\bottomrule
\end{tabular}
\end{table}

\begin{important}
这组接口是"新增 PDE 只需十几行"的根本原因：
物理层永远只做矩阵的线性组合与边界处理，不必关心单元循环。
\end{important}

\section{Lagrange 元（经典有限元）}

\subsection{P1 三角形}

参考三角形上，用重心坐标 $\lambda_0=1-\xi-\eta$，$\lambda_1=\xi$，$\lambda_2=\eta$：

\begin{equation}
\hat\varphi_0=1-\xi-\eta,\qquad
\hat\varphi_1=\xi,\qquad
\hat\varphi_2=\eta .
\label{eq:p1}
\end{equation}

\subsection{P2 三角形}

\begin{align}
\hat\varphi_0&=\lambda_0(2\lambda_0-1),&
\hat\varphi_1&=\lambda_1(2\lambda_1-1),&
\hat\varphi_2&=\lambda_2(2\lambda_2-1),\notag\\
\hat\varphi_3&=4\lambda_0\lambda_1\ (\text{棱 }0\text{-}1),&
\hat\varphi_4&=4\lambda_1\lambda_2\ (\text{棱 }1\text{-}2),&
\hat\varphi_5&=4\lambda_2\lambda_0\ (\text{棱 }2\text{-}0).
\label{eq:p2}
\end{align}

\subsection{张量积元}

四边形与六面体使用 1D Lagrange 基的张量积。以 Q2 为例，
1D 节点取 $\{0,\tfrac12,1\}$：

\begin{equation}
N_0(t)=2t^2-3t+1,\qquad
N_1(t)=-4t^2+4t,\qquad
N_2(t)=2t^2-t,
\label{eq:q2-1d}
\end{equation}

单元基函数为 $N_{ijk}(\xi,\eta,\zeta)=N_i(\xi)N_j(\eta)N_k(\zeta)$。
由于节点集是完整张量网格，
$\{N_{ijk}\}$ 自动构成对偶于节点值的\textbf{节点基}。

\subsection{自由度的全局编号}

\mosica{} 用一个"登记簿"把几何实体映射到全局自由度：

\begin{itemize}[itemsep=0.15em]
  \item 顶点自由度：以全局顶点编号为 key；
  \item 棱自由度：以两个端点的\textbf{有序对}为 key；
  \item 面自由度（3D）：以面的 4 个顶点有序元组为 key；
  \item 单元内部自由度：以单元编号为 key。
\end{itemize}

于是相邻单元自动共享棱 / 面自由度，全局空间是 $H^1$ 协调的。

\begin{lstlisting}[language=MosicaFE,caption={示例 5.1：创建空间并查看自由度}]
from MosicaFE import Mesh, FESpace

mesh = Mesh.rectangle(nx=8, ny=8)

V1 = FESpace.lagrange(mesh, degree=1)
V2 = FESpace.lagrange(mesh, degree=2)
print(V1.n_dofs, V2.n_dofs)      # 81 289

print(V2.get_cell_dofs(0))       # local -> global DOF map
print(V2.get_dof_coordinates(0))
print(V2.boundary_dofs()[:6])
\end{lstlisting}

\section{虚拟元空间}

\begin{lstlisting}[language=MosicaFE,caption={示例 5.2：虚拟元空间（任意多边形）}]
from MosicaFE import Mesh, FESpace

mesh = Mesh.voronoi_polygon(n_points=100, seed=0)
V = FESpace.vem(mesh)                       # k = 1, vertex DOFs only

print(V.n_dofs, mesh.n_nodes)               # equal
K = V.stiffness_matrix()                    # projection + stabilization
M = V.mass_matrix()
b = V.load_vector(lambda x: 1.0)

print("row sums of M equal b?",
      abs(b - M.sum(axis=1)).max() < 1e-12)
\end{lstlisting}

\begin{property}
\code{VEMSpace} 的刚度矩阵满足
\begin{enumerate}[label=(\arabic*),itemsep=0.1em]
  \item 对称（因为稳定化项写成 Gram 矩阵的形式）；
  \item 半正定，核空间中包含常函数；
  \item 一次多项式被精确再生（单元矩阵作用在任意线性自由度向量上，
        给出的能量等于 $\int_K|\nabla p|^2$）。
\end{enumerate}
第 (3) 条就是"patch test"，第 \ref{sec:pitfall-vem} 节会说明它有多重要。
\end{property}

\begin{table}[htbp]
\centering
\caption{\code{VEMSpace} 的 \code{stabilization} 选项}
\label{tab:stab}
\begin{tabular}{@{}lp{8.6cm}@{}}
\toprule
取值 & 含义 \\
\midrule
\code{"trace"}（默认） & $\sigma(I-\Pi_{\rm dof})^{\top}(I-\Pi_{\rm dof})$，
对称半正定，$\sigma$ 由迹比确定。推荐使用 \\
\code{"legacy"} & 早期脚本的 $(I-\Pi_{\rm dof})/n$，非对称，
在不规则多边形上会发散，仅用于复现旧结果 \\
\code{"none"} & 只用投影项 $A_{\rm proj}$。单元矩阵是奇异的，
仅用于诊断，不要用于求解 \\
\bottomrule
\end{tabular}
\end{table}

\section{RT0$\times$P0 混合空间}

\begin{lstlisting}[language=MosicaFE,caption={示例 5.3：混合空间}]
from MosicaFE import Mesh, FESpace

mesh = Mesh.rectangle(nx=16, ny=16)       # triangle mesh
V = FESpace.rt0(mesh)

print(V.n_flux, V.n_pressure, V.n_dofs)   # facets, cells, sum
print(V.cell_flux_dofs(0))                # local facet -> global flux DOF
print(V.facet_signs(0))                   # outward orientation signs
print(V.local_flux_mass(0))               # 单元 0 的精确 RT0 质量矩阵 (3x3)
\end{lstlisting}

\subsection{空间层给出的两个块}

RT0 空间与 \code{LagrangeSpace} / \code{VEMSpace} 的接口是\textbf{平行}的，
只是把"单个椭圆算子"换成了"鞍点问题的两个块"：

\begin{table}[htbp]
\centering
\caption{RT0$\times$P0 的矩阵原语（供物理层组合）}
\label{tab:rt0-blocks}
\begin{tabular}{@{}L{5.4cm}L{2.9cm}L{6.1cm}@{}}
\toprule
方法 & 形状 & 含义 \\
\midrule
\code{V.flux\_mass\_matrix()} & $(\Phi,\Phi)$ & $M_{ij}=\int_K\varphi_i\cdot\varphi_j$，\textbf{精确闭式}，无需求积 \\
\code{V.divergence\_matrix()} & $(P,\Phi)$ & $B_{K,i}=\int_K\nabla\cdot\varphi_i=\sigma_{K,i}$，恒为 $\pm1$ \\
\code{V.pressure\_load\_vector(f)} & $(P,)$ & $\int_K f$ \\
\midrule
\code{V.interpolate\_pressure(f)} & $(P,)$ & $f$ 的单元平均（$L^2$ 投影到 P0） \\
\code{V.interpolate\_flux(q)} & $(\Phi,)$ & $\Phi_f=\int_f q\cdot n_f\dif s$ \\
\code{V.evaluate\_flux(c,Q,pts)} & $(n_q,d)$ & 单元内重构 $q_h=\sum_i F_i\varphi_i$ \\
\code{V.locate\_cell(x)} & int & 用重心坐标定位点所在单元 \\
\bottomrule
\end{tabular}
\end{table}

\begin{lstlisting}[language=MosicaFE,caption={示例 5.4：直接使用两个块}]
from MosicaFE import Mesh, FESpace
from MosicaFE._compat import bmat

mesh = Mesh.rectangle(nx=16, ny=16)
V = FESpace.rt0(mesh)

M = V.flux_mass_matrix()          # (Phi,Phi)
B = V.divergence_matrix()         # (P,Phi)
f = V.pressure_load_vector(src)   # (P,)

A = bmat([[M, -B.T], [B, None]])  # [[M, -B^T], [B, 0]]
\end{lstlisting}

\begin{caution}
RT0 空间\textbf{不提供} \code{stiffness\_matrix()} / \code{mass\_matrix()} /
\code{load\_vector()}，也不提供 \code{interpolate()}：它对应的是鞍点问题，
而且同时携带两个场，没有"单一算子 / 单一插值"可言。
调用这些方法会抛出 \code{NotImplementedError}，并在消息里指出该用哪个方法
（例如 \code{mass\_matrix()} $\to$ \code{flux\_mass\_matrix()}，
\code{load\_vector()} $\to$ \code{pressure\_load\_vector()}）。
日常使用请直接配合 \code{MixedPoissonProblem}。
\end{caution}

\section{空间工厂 \code{FESpace}}

\begin{lstlisting}[language=MosicaFE,caption={示例 5.4：三种空间族}]
from MosicaFE import Mesh, FESpace

mesh = Mesh.rectangle(nx=16, ny=16)

V_fem = FESpace.lagrange(mesh, degree=2)    # P2
V_vem = FESpace.vem(mesh, degree=1)         # VEM k=1
V_rt0 = FESpace.rt0(mesh)                   # RT0 x P0

# by name (used by solve_poisson(element=...))
print(FESpace.available())
V_any = FESpace.create(mesh, family="vem", degree=1)
\end{lstlisting}

\begin{table}[htbp]
\centering
\caption{字符串形式的"元"名字（\code{solve\_poisson(element=...)} 使用）}
\label{tab:elem-names}
\begin{tabular}{@{}L{3.6cm}L{3.6cm}L{6.2cm}@{}}
\toprule
名字 & 空间族 & 可用网格 \\
\midrule
\code{"P1"} / \code{"P2"} & Lagrange & 三角形、四面体 \\
\code{"Q1"} / \code{"Q2"} & Lagrange（张量积） & 四边形、六面体 \\
\code{"VEM"} / \code{"VEM1"} & 虚拟元 $k=1$ & 任意多边形 / 多面体 \\
\code{"RT0"} / \code{"RT0-P0"} / \code{"MIXED"} & RT0$\times$P0 & 三角形、四面体 \\
\bottomrule
\end{tabular}
\end{table}

% ===========================================================================
\chapter{物理问题与组装}
\label{chap:physics}

\section{问题接口}

所有 PDE 都继承 \code{BasePhysics}，只需实现一个方法：

\begin{lstlisting}[language=MosicaFE,caption={示例 6.1：\code{BasePhysics} 接口},label={lst:basephysics}]
class BasePhysics(ABC):
    def __init__(self, space, f=None, g=None, **kwargs):
        self.space = space
        self.mesh = space.mesh
        self.n_dofs = space.n_dofs
        self.f = f or (lambda x: 0.0)      # source term
        self.g = g or (lambda x: 0.0)      # Dirichlet data
        self.boundary_dofs = space.boundary_dofs()

    @abstractmethod
    def assemble(self):
        """return (A, b) with boundary conditions applied"""

    def solve(self, method="auto", **kwargs):
        A, b = self.assemble()
        return Solver.solve(A, b, method=method, **kwargs)
\end{lstlisting}

\section{Poisson 问题}

\begin{lstlisting}[language=MosicaFE,caption={示例 6.2：Poisson 问题的组装（完整实现）}]
@register_problem("poisson")
class PoissonProblem(BasePhysics):
    def __init__(self, space, f=None, g=None, kappa=1.0, **kwargs):
        super().__init__(space, f=f, g=g, **kwargs)
        self.kappa = kappa

    def assemble(self):
        A = self.space.stiffness_matrix()
        if float(self.kappa) != 1.0:
            A = float(self.kappa) * A
        b = self.space.load_vector(self.f)
        A, b = self.apply_boundary_conditions(A, b)
        return A, b
\end{lstlisting}

\begin{important}
整个实现只有 8 行，而且它对 \textbf{Lagrange 元与虚拟元同时成立}。
这正是把"单元循环"下沉到 \module{spaces} 层之后的收益。
\end{important}

\section{边界条件}

\subsection{强加 Dirichlet 条件的正确做法}

对边界自由度 $j$ 强加 $u_j=g_j$，标准做法是"消行消列 + 右端项修正"：

\begin{enumerate}[label=(\arabic*),itemsep=0.15em]
  \item 把边界列对内部行的贡献搬到右端项：
        \begin{equation}
        b_i \leftarrow b_i-\sum_{j\in\mathcal{D}}A_{ij}g_j,
        \qquad \forall i\notin\mathcal{D};
        \label{eq:bc-lift}
        \end{equation}
  \item 删去所有涉及 $\mathcal{D}$ 的行与列；
  \item 对 $j\in\mathcal{D}$ 置 $A_{jj}=1$，$b_j=g_j$。
\end{enumerate}

\begin{caution}
第 (1) 步在齐次条件 $g=0$ 时"看起来"可以省略，
因此很多演示程序都漏掉了它；一旦边界数据非零，解就是错的。
\mosica{} 用一个简单而灵敏的测试——\textbf{非齐次线性 patch test}——
来防止这类错误回归：

\begin{equation}
u(x)=1+2x_1-\tfrac12 x_2,\qquad f\equiv0,\qquad g=u|_{\partial\Omega},
\label{eq:patch}
\end{equation}

此时精确解属于离散空间，数值解应当\textbf{精确}（误差 $\sim10^{-15}$）。
修复前的实测误差是 $2.875$。
\end{caution}

\subsection{其它边界条件}

\begin{itemize}[itemsep=0.2em]
  \item \textbf{Neumann 条件}：向 \code{load\_vector} 的结果中加入
        $\int_{\Gamma_N}h\varphi_i\dif s$ 即可，需要边界 facet 的求积规则；
  \item \textbf{RT0 混合格式}：$p=0$ 是自然边界条件，
        无需消去任何自由度（这也是混合格式在流量问题上更"自然"的原因）；
  \item 非齐次 $p=g\ne0$ 在混合格式下需要额外处理边界通量项，
        当前版本会显式报错而不是给出错误结果。
\end{itemize}

\section{组装算法}

\begin{algorithm}[htbp]
\caption{全局组装（\mosica{} 的 \code{BaseFESpace.\_assemble}）}
\label{alg:assembly}
\begin{algorithmic}[1]
\REQUIRE 空间 $V$、单元局部矩阵生成函数 \code{local\_matrix(cell)}
\STATE 初始化三元组列表 \code{rows, cols, vals}
\FOR{每个单元 $c=0,\dots,n_{\rm cells}-1$}
    \STATE $\mathcal{D}_c\leftarrow$\code{get\_cell\_dofs(c)}
    \STATE $K_c\leftarrow$\code{local\_matrix(c)}
    \STATE 把 $\mathcal{D}_c$ 的行 / 列索引与 $K_c$ 展平后追加到三元组列表
\ENDFOR
\STATE 由 COO 三元组构造矩阵（重复下标自动求和）
\IF{要求对称}
    \STATE $A\leftarrow\frac12(A+A^{\top})$
\ENDIF
\RETURN $A$
\end{algorithmic}
\end{algorithm}

\begin{caution}
右端项组装必须使用"逐项累加"语义（\code{numpy} 的 \code{np.add.at}）。
若直接用 \code{b[dofs] += local}，当 \code{dofs} 中出现重复下标时
（例如网格里存在重复顶点），\code{numpy} 的缓冲赋值会丢掉累加，
而稀疏矩阵的 COO 组装却会正确求和——两者不一致会造成难以察觉的错误。
\end{caution}

\section{使用 Poisson 问题}

\begin{lstlisting}[language=MosicaFE,caption={示例 6.3：组装与求解}]
import numpy as np
from MosicaFE import Mesh, FESpace, PoissonProblem

def source(x):
    return 2*np.pi**2 * np.sin(np.pi*x[0]) * np.sin(np.pi*x[1])

mesh = Mesh.rectangle(nx=32, ny=32)
V = FESpace.lagrange(mesh, degree=2)

problem = PoissonProblem(V, f=source, g=lambda x: 0.0, kappa=1.0)

A, b = problem.assemble()
print("system size:", A.shape, "non-zeros:", A.nnz)

u = problem.solve(method="direct")
print("solution size:", u.size)
\end{lstlisting}

\section{PDE 注册表}

\begin{lstlisting}[language=MosicaFE,caption={示例 6.4：按名字创建与求解方程}]
from MosicaFE import FESpace, Mesh, available_problems, create_problem, run

print(available_problems())        # ['mixed_poisson', 'poisson', ...]

mesh = Mesh.rectangle(nx=32, ny=32)
problem = create_problem("poisson", FESpace.vem(mesh),
                         f=source, g=lambda x: 0.0)
u = problem.solve()

u2 = run("poisson", mesh, element="VEM", f=source, g=lambda x: 0.0)
\end{lstlisting}

% ===========================================================================
\chapter{线性求解器}
\label{chap:solvers}

\section{求解器接口}

\begin{lstlisting}[language=MosicaFE,caption={示例 7.1：\code{Solver} 接口}]
import numpy as np, time
from MosicaFE import FESpace, Mesh, PoissonProblem, Solver

mesh = Mesh.rectangle(nx=32, ny=32)
problem = PoissonProblem(FESpace.lagrange(mesh, 2), f=source,
                         g=lambda x: 0.0)
A, b = problem.assemble()

t0 = time.time(); u_direct = Solver.direct(A, b); t1 = time.time()
u_cg    = Solver.conjugate_gradient(A, b, tol=1e-10)
u_gmres = Solver.gmres(A, b, tol=1e-10, restart=30)
u_bicg  = Solver.bicgstab(A, b, tol=1e-10)

print("direct %.3f s, diff vs CG: %.3e"
      % (t1 - t0, np.linalg.norm(u_direct - u_cg)))
\end{lstlisting}

也可以让问题对象自己选择：

\begin{lstlisting}[language=MosicaFE]
u = problem.solve(method="auto")   # auto/direct/cg/gmres/bicgstab/minres
\end{lstlisting}

\section{直接法}

\begin{itemize}[itemsep=0.2em]
  \item 稀疏矩阵：\code{scipy.sparse.linalg.spsolve}（内部为稀疏 LU / Cholesky）；
  \item 稠密矩阵：\code{numpy.linalg.solve}；
  \item 复杂度：2D 约 $O(n^{1.5})$，3D 约 $O(n^{2})$。
\end{itemize}

\section{迭代法}

\begin{table}[htbp]
\centering
\caption{可用的 Krylov 求解器}
\label{tab:krylov}
\begin{tabular}{@{}llp{6.4cm}@{}}
\toprule
方法 & 适用矩阵 & 备注 \\
\midrule
\code{conjugate\_gradient} & 对称正定 & 每步 $O(\mathrm{nnz})$，迭代数 $O(\sqrt{\kappa})$ \\
\code{gmres} & 一般（非对称） & 带重启，存储 $O(\text{restart}\times n)$ \\
\code{bicgstab} & 一般 & 每步两次矩阵向量积，内存小 \\
\code{minres} & 对称不定 & 例如 RT0 的鞍点系统 \\
\bottomrule
\end{tabular}
\end{table}

\begin{caution}
共轭梯度法要求矩阵对称正定。虚拟元的刚度矩阵在默认稳定化下是对称的，
但如果选了 \code{stabilization="legacy"}，矩阵不再对称，CG 会失效
（甚至不报错地给出错误结果）。鞍点系统请用直接法或 MINRES。
\end{caution}

\section{求解器选择指南}

\begin{table}[htbp]
\centering
\caption{求解器选择建议}
\label{tab:solver-guide}
\begin{tabular}{@{}llll@{}}
\toprule
问题规模 & 矩阵类型 & 推荐求解器 & 备注 \\
\midrule
$<2\times10^4$ 自由度 & 任意 & 直接法 & 精确、无需调参 \\
$>2\times10^4$ & 对称正定 & CG & 建议配预处理 \\
$>2\times10^4$ & 对称不定 & MINRES 或直接法 & 鞍点问题 \\
$>2\times10^4$ & 非对称 & GMRES / BiCGSTAB & 需要重启参数 \\
\bottomrule
\end{tabular}
\end{table}

\begin{lstlisting}[language=MosicaFE,caption={示例 7.2：自动选择}]
from MosicaFE import Solver

print(Solver.recommend(1000, symmetric=True))     # 'direct'
print(Solver.recommend(10**6, symmetric=True))    # 'cg'
print(Solver.recommend(10**6, symmetric=False))   # 'gmres'
\end{lstlisting}

% ===========================================================================
\chapter{误差分析与收敛性}
\label{chap:error}

\section{误差范数}

\begin{align}
\|e\|_{L^2}&=\Bigl(\int_\Omega|u-u_h|^2\dif x\Bigr)^{1/2},&
\|e\|_{L^2}^2&\approx\sum_K\sum_q w_q\,|u(x_q)-u_h(x_q)|^2|\det J_K|,
\label{eq:l2norm}\\
|e|_{H^1}&=\Bigl(\int_\Omega|\nabla u-\nabla u_h|^2\dif x\Bigr)^{1/2},&
\|e\|_{L^\infty}&=\max_{x\in\Omega}|u(x)-u_h(x)| .
\label{eq:h1norm}
\end{align}

\begin{lstlisting}[language=MosicaFE,caption={示例 8.1：计算三种范数}]
import numpy as np

def exact(x):
    return np.sin(np.pi*x[0]) * np.sin(np.pi*x[1])

def exact_grad(x):
    return np.array([np.pi*np.cos(np.pi*x[0])*np.sin(np.pi*x[1]),
                     np.pi*np.sin(np.pi*x[0])*np.cos(np.pi*x[1])])

u = problem.solve()
l2   = problem.compute_error(u, exact, "L2")
linf = problem.compute_error(u, exact, "Linf")
h1   = problem.compute_error(u, exact, "H1", exact_grad=exact_grad)

print("%.6e %.6e %.6e" % (l2, h1, linf))

norms = problem.compute_errors(u, exact, exact_grad)
\end{lstlisting}

\begin{remark}
虚拟元的基函数没有显式表达式，\code{VEMSpace.evaluate} 用的是
"顶点自由度上的线性最小二乘重构"。该重构误差与 $u_h$ 本身同阶，
因此实测收敛阶与 $u_h$ 的真实收敛阶一致（见第 \ref{sec:verify} 节）。
\end{remark}

\section{收敛性研究框架}

\begin{lstlisting}[language=MosicaFE,caption={示例 8.2：自动收敛性研究},label={lst:convstudy}]
from MosicaFE import FESpace, Mesh, PoissonProblem, convergence_study
from MosicaFE.postprocess import estimate_convergence_rate

def mesh_factory(level):
    n = 4 * 2**level
    return Mesh.rectangle(nx=n, ny=n)

def problem_factory(mesh):
    V = FESpace.lagrange(mesh, degree=2)
    return PoissonProblem(V, f=source, g=lambda x: 0.0)

h, errors = convergence_study(mesh_factory, problem_factory,
                              exact, exact_grad,
                              refinement_levels=5, plot=True,
                              filename="conv_p2.png")

print(estimate_convergence_rate(h, errors["L2"], use_last=3))   # ~3.0
\end{lstlisting}

\section{收敛阶估计}

由两层级误差

\begin{equation}
\text{rate}=\frac{\log(e_1/e_2)}{\log(h_1/h_2)} .
\label{eq:rate}
\end{equation}

更稳健的做法是对 $\log e$--$\log h$ 做最小二乘拟合
（\code{estimate\_convergence\_rate(h, errors)} 就是这样实现的），
并且可以用 \code{use\_last=k} 只取最后 $k$ 层以逼近渐近阶。

\section{实测验证结果}
\label{sec:verify}

\begin{table}[htbp]
\centering
\caption{各方法的实测收敛阶（制造解，网格逐次加密）}
\label{tab:verify}
\begin{tabular}{@{}llccc@{}}
\toprule
元 & 网格 & 实测 $L^2$ & 实测 $H^1$ & 理论 $(L^2/H^1)$ \\
\midrule
P1 & 三角形 & 1.98 & 0.99 & 2 / 1 \\
P2 & 三角形 & 3.00 & 1.99 & 3 / 2 \\
Q1 & 四边形 & 2.00 & 1.00 & 2 / 1 \\
Q2 & 四边形 & 3.00 & 2.00 & 3 / 2 \\
VEM $k=1$ & 三角形 & 1.98 & 0.99 & 2 / 1 \\
VEM $k=1$ & 四边形 & 1.99 & 1.00 & 2 / 1 \\
VEM $k=1$ & 五边形 & 2.00 & 1.00 & 2 / 1 \\
VEM $k=1$ & Voronoi 多边形 & 1.97 & 1.01 & 2 / 1 \\
P1 & 四面体 & 1.95 & 0.98 & 2 / 1 \\
Q1 & 六面体 & 2.01 & 1.00 & 2 / 1 \\
VEM $k=1$ & 六面体 & 1.97 & 0.97 & 2 / 1 \\
RT0$\times$P0 & 三角形（压力） & 1.00 & --- & 1 \\
RT0$\times$P0 & 三角形（通量） & 1.00 & --- & 1 \\
\bottomrule
\end{tabular}
\end{table}

复现方式：

\begin{lstlisting}[style=shell]
& "D:\anaconda\envs\fealpy\python.exe" MosicaFE\tests\run_checks.py
\end{lstlisting}

% ===========================================================================
\chapter{可视化}
\label{chap:visual}

\section{可视化工具一览}

\begin{table}[htbp]
\centering
\caption{\module{postprocess} 中的绘图函数}
\label{tab:plots}
\begin{tabular}{@{}L{5.6cm}L{9.8cm}@{}}
\toprule
函数 & 用途 \\
\midrule
\code{plot\_solution(u, mesh, V, ...)} & 画解。1D 画曲线；2D 画等值线填充或三维曲面；
3D 画 $z=\text{const}$ 切片。多边形网格自动从形心做扇形三角化 \\
\code{plot\_mesh(mesh, ...)} & 画网格（3D 画中间切片） \\
\code{plot\_error(u, mesh, space, exact, ...)} & 画逐点误差 $|u_h-u|$ \\
\code{plot\_convergence(h, errors, ...)} & log--log 收敛曲线，并叠加 $O(h^k)$ 参考斜率 \\
\code{plot\_pressure\_2d(mesh, P, ...)} & RT0 的分片常数压力 \\
\code{plot\_flux\_2d(mesh, Q, ...)} & 通量箭头图 \\
\code{compare\_solutions(entries, ...)} & 多组解并排比较 \\
\bottomrule
\end{tabular}
\end{table}

通用参数：\code{title=}、\code{filename=}（保存 PNG）、\code{show=}（弹窗）、
\code{dpi=}、\code{ax=}（画到已有坐标轴上以便拼图）。

\section{画解}

\begin{lstlisting}[language=MosicaFE,caption={示例 9.1：2D 与 3D 解的绘制}]
from MosicaFE import Mesh, FESpace, PoissonProblem
from MosicaFE.postprocess import plot_solution, plot_mesh, plot_error

# --- 2D contour ---
mesh = Mesh.rectangle(nx=32, ny=32)
V = FESpace.lagrange(mesh, 2)
u = PoissonProblem(V, f=source, g=lambda x: 0.0).solve()
plot_solution(u, mesh, V, title="FEM P2", filename="sol2d.png")

# --- 2D surface ---
plot_solution(u, mesh, V, mode="surface", filename="sol_surface.png")

# --- polygon mesh + VEM solution ---
mp = Mesh.voronoi_polygon(n_points=200, seed=3)
Vv = FESpace.vem(mp)
uv = PoissonProblem(Vv, f=source, g=lambda x: 0.0).solve()
plot_solution(uv, mp, Vv, title="VEM on polygons",
              filename="sol_vem.png")
plot_error(uv, mp, Vv, exact, filename="err_vem.png")

# --- 3D slice ---
m3 = Mesh.box(nx=8, ny=8, nz=8)
V3 = FESpace.lagrange(m3, 1)
u3 = PoissonProblem(V3, f=source3d, g=lambda x: 0.0).solve()
plot_solution(u3, m3, V3, z_slice=0.5, filename="sol3d_slice.png")
plot_mesh(m3, filename="mesh3d.png")
\end{lstlisting}

\section{画收敛曲线}

\begin{lstlisting}[language=MosicaFE,caption={示例 9.2：收敛曲线}]
from MosicaFE.postprocess import plot_convergence

plot_convergence(h, errors, filename="convergence.png",
                 reference_rates=(1, 2, 3))
\end{lstlisting}

\begin{caution}
所有绘图函数在没有安装 \code{matplotlib} 时只会打印提示并返回 \code{None}，
不会抛异常，因此批量计算的脚本可以放心调用。
\end{caution}

% ===========================================================================
\chapter{完整示例}
\label{chap:examples}

\section{示例一：两行求解 Poisson 方程}

\begin{lstlisting}[language=MosicaFE,caption={示例 10.1：一行式求解}]
import numpy as np
from MosicaFE import Mesh, solve_poisson

source = lambda x: 2*np.pi**2 * np.sin(np.pi*x[0]) * np.sin(np.pi*x[1])

u = solve_poisson(Mesh.rectangle(nx=32, ny=32), element="P2", f=source)
print(u.size)          # 4225
\end{lstlisting}

\section{示例二：换一种元（同一段代码，三种方法）}

\begin{lstlisting}[language=MosicaFE,caption={示例 10.2：FEM / VEM / RT0 共用同一套流程},label={lst:three}]
import numpy as np
from MosicaFE import FESpace, Mesh, MixedPoissonProblem, PoissonProblem

def exact(x):
    return np.sin(np.pi*x[0]) * np.sin(np.pi*x[1])

def grad(x):
    return np.array([np.pi*np.cos(np.pi*x[0])*np.sin(np.pi*x[1]),
                     np.pi*np.sin(np.pi*x[0])*np.cos(np.pi*x[1])])

def source(x):
    return 2*np.pi**2 * exact(x)

def zero(x):
    return 0.0

mesh = Mesh.rectangle(nx=32, ny=32)

# (a) classic FEM: P2 Lagrange
V = FESpace.lagrange(mesh, degree=2)
p = PoissonProblem(V, f=source, g=zero)
u = p.solve()
print("FEM P2 :", p.compute_errors(u, exact, grad))

# (b) VEM: only the space changes
V = FESpace.vem(mesh)
p = PoissonProblem(V, f=source, g=zero)
u = p.solve()
print("VEM k=1:", p.compute_errors(u, exact, grad))

# (c) RT0 x P0 mixed: space + problem class
V = FESpace.rt0(mesh)
p = MixedPoissonProblem(V, f=source)
Q, P = p.solve()
print("RT0-P0 :", p.compute_errors((Q, P), exact, lambda x: -grad(x)))
\end{lstlisting}

运行结果（32$\times$32 网格）：

\begin{lstlisting}[style=textout]
FEM P2 : {'L2': 6.87e-05, 'Linf': 2.41e-03, 'H1': 7.38e-03}
VEM k=1: {'L2': 1.35e-03, 'Linf': 3.18e-03, 'H1': 1.09e-01}
RT0-P0 : {'L2': 1.64e-02, 'Linf': 4.55e-02, 'flux_L2': 6.30e-02}
\end{lstlisting}

\section{示例三：多边形网格上的虚拟元}

\begin{lstlisting}[language=MosicaFE,caption={示例 10.3：Voronoi 多边形网格 + 收敛性研究}]
import numpy as np
from MosicaFE import FESpace, Mesh, PoissonProblem, convergence_study
from MosicaFE.postprocess import estimate_convergence_rate, plot_solution

def mesh_factory(level):
    return Mesh.voronoi_polygon(n_points=50 * 4**level, seed=7)

def problem_factory(mesh):
    return PoissonProblem(FESpace.vem(mesh), f=source, g=lambda x: 0.0)

h, errors = convergence_study(mesh_factory, problem_factory,
                              exact, grad, refinement_levels=3,
                              verbose=True)
print("L2 rate:", estimate_convergence_rate(h, errors["L2"]))

mesh = Mesh.voronoi_polygon(n_points=200, seed=7)
V = FESpace.vem(mesh)
u = PoissonProblem(V, f=source, g=lambda x: 0.0).solve()
plot_solution(u, mesh, V, title="VEM on Voronoi mesh",
              filename="vem_voronoi.png")
\end{lstlisting}

\begin{caution}
随机 Voronoi 网格的 $h$ 用"最大单元直径"会很不稳定，
建议改用平均单元直径作为横坐标，\code{convergence\_study} 里传自定义
\code{mesh\_factory} 即可（示例中用的是 \code{mesh.get\_mesh\_size()}，
若要更稳健可在回调里替换）。
\end{caution}

\section{示例四：自定义新 PDE（Helmholtz）}

\begin{lstlisting}[language=MosicaFE,caption={示例 10.4：新增一门方程只需十几行},label={lst:custompde}]
import numpy as np
from MosicaFE import FESpace, Mesh
from MosicaFE.physics.base import BasePhysics
from MosicaFE.physics.registry import register_problem

@register_problem("helmholtz")
class HelmholtzProblem(BasePhysics):
    """-Laplacian(u) - k^2 u = f, weak form: (grad u, grad v) - k^2 (u, v) = (f, v)"""

    def __init__(self, space, f=None, g=None, k=1.0, **kwargs):
        super().__init__(space, f=f, g=g, **kwargs)
        self.k = k

    def assemble(self):
        A = self.space.stiffness_matrix() - self.k**2 * self.space.mass_matrix()
        b = self.space.load_vector(self.f)
        return self.apply_boundary_conditions(A, b)

# use it with any space family
k = 1.0
exact = lambda x: np.sin(np.pi*x[0]) * np.sin(np.pi*x[1])
source = lambda x: (2*np.pi**2 - k**2) * exact(x)

mesh = Mesh.rectangle(nx=32, ny=32)
V = FESpace.vem(mesh)                       # works with VEM too
problem = HelmholtzProblem(V, f=source, g=lambda x: 0.0, k=k)
u = problem.solve(method="direct")
print("max error:", np.abs(u - V.interpolate(exact)).max())
\end{lstlisting}

输出：\code{max error: 9.310e-04}。

\section{示例五：反应扩散方程}

\begin{lstlisting}[language=MosicaFE,caption={示例 10.5：内置的反应扩散问题}]
import numpy as np
from MosicaFE import FESpace, Mesh, ReactionDiffusionProblem

c = 10.0
exact = lambda x: np.sin(np.pi*x[0]) * np.sin(np.pi*x[1])
source = lambda x: (2*np.pi**2 + c) * exact(x)

mesh = Mesh.rectangle(nx=32, ny=32)
V = FESpace.lagrange(mesh, degree=2)
problem = ReactionDiffusionProblem(V, f=source, g=lambda x: 0.0,
                                   reaction_coeff=c)
u = problem.solve()
print(problem.compute_error(u, exact, "L2"))     # ~6.8e-05
\end{lstlisting}

% ===========================================================================
\chapter{进阶话题}
\label{chap:advanced}

\section{精度陷阱与经验教训}
\label{sec:pitfall-vem}

把三套程序整合成一个库的过程中，实测暴露了四类问题。
它们都不是"数值精度的小差别"，而是会给出错误结果或让收敛阶崩掉的结构性错误。
记录在此，既是为了说明 \mosica{} 为什么这么写，也供读者在自己的代码里自查。

\subsection{陷阱一：Dirichlet 消去漏掉右端项修正}

\textbf{症状}：齐次边界一切正常，非齐次边界解错。

\textbf{诊断}：用线性 patch test（式 \eqref{eq:patch}）。
精确解属于离散空间时，数值解必须精确。

\textbf{结论}：见式 \eqref{eq:bc-lift}，第 (1) 步不可省。
\mosica{} 把这条测试写进了 \code{tests/test\_core\_spaces.py}。

\subsection{陷阱二：张量积单元的求积阶不足}

\textbf{症状}：三维六面体 Q1 的 $L^2$ 收敛阶算出来是 $-0.11$。

\textbf{诊断}：Q1 刚度矩阵的被积函数在每个方向上是 2 次多项式，
而每个方向只用了 1 个 Gauss 点时只能精确到 1 次，欠积分导致
单元矩阵不是正定的、能量被系统性低估。

\textbf{结论}：张量积单元的求积阶取 \code{max(2, 2*p)}（每方向 $p$ 个点起步），
见表格 \ref{tab:quad} 后的注意事项。

\subsection{陷阱三：虚拟元投影的常数模式规范化}

\textbf{症状}：三角形与正方形网格上 VEM 收敛阶正常，
但换到 Voronoi 多边形网格，$L^2$ 收敛阶从 2 掉到 0.5 以下，
误差在 $10^{-2}$ 处停滞。

\textbf{诊断}：把线性精确解代入，检查数值解是否精确（patch test），
发现不规则多边形上失败。

\textbf{结论}：常数模式必须按式 \eqref{eq:vem-const} 规范化，
使投影精确再生一次多项式；否则稳定化项会惩罚线性函数，
破坏相容性。

\subsection{陷阱四：载荷 / 质量矩阵的几何权重}

\textbf{症状}：即使修好了投影，不规则多边形网格上的 $L^2$ 误差仍然偏大。

\textbf{诊断}：把 $\int_K\varphi_i$ 一律取成 $|K|/n$，
相当于假设每个顶点"平分"单元。对三角形、正四边形这是对的，
但对边长差异很大的多边形，短边上的顶点会被严重高估。

\textbf{结论}：改用式 \eqref{eq:vem-load} 的扇形线性插值，
几何权重由单元自身形状决定。修复后 Voronoi 网格的
$L^2$ 收敛阶回到 1.97。

\subsection{陷阱五：numpy 的重复下标赋值}

\textbf{症状}：\code{b[dofs] += local} 与稀疏矩阵组装结果不一致，
质量矩阵的行和与载荷向量对不上。

\textbf{诊断}：\code{dofs} 中出现重复下标时
（例如网格里有重复顶点），\code{numpy} 的缓冲赋值只保留最后一次，
而 COO 组装会求和。

\textbf{结论}：一律使用 \code{numpy.add.at}。

\section{网格细化策略}

\begin{itemize}[itemsep=0.2em]
  \item \textbf{均匀加密}：\code{mesh.refine\_uniform()}，
        三角形单元数 $\times4$、六面体 $\times8$；
  \item \textbf{自适应加密（规划中）}：基于残差型误差指示子
        \begin{equation}
        \eta_K=h_K\|f+\Delta u_h\|_{L^2(K)}
        +h_K^{1/2}\|[\partial_n u_h]\|_{L^2(\partial K)}
        \label{eq:estimator}
        \end{equation}
        标记 $\eta_K>\theta\max_K\eta_K$ 的单元细分。
        \mosica{} 的 \code{Mesh} 支持变长单元，
        因此局部加密在数据结构上是可行的，缺的是 refinement 的实现。
\end{itemize}

\section{性能与向量化}

\begin{enumerate}[label=(\arabic*),itemsep=0.2em]
  \item 组装阶段累计 COO 三元组，最后一次转成 CSR；
  \item 参考单元上的基函数求值按"求积点"缓存
        （\code{LagrangeSpace.\_quadrature\_data}），避免逐单元重复计算；
  \item 对同一网格上多种物理问题，\code{stiffness\_matrix()} 的结果可以复用；
  \item 大规模问题建议改用迭代法；
  \item 若要进一步提速，可以把单元循环用 \code{numba} 或
        "按单元类型分块 + 批量张量运算"的方式向量化。
\end{enumerate}

\section{与其它库的关系}

\mosica{} 刻意保持"小、可读、可改"：

\begin{itemize}[itemsep=0.2em]
  \item 与 FEniCS / Firedrake 相比，\mosica{} 不是通用有限元平台，
        而是\textbf{教学与研究用的可读实现}，代价是只支持结构化网格生成
        与有限几类单元；
  \item 与 \code{fealpy} 相比，\mosica{} 更聚焦于
        "FEM / VEM / RT0 三种离散方式的横向对比"；
  \item 与原始三套脚本相比，\mosica{} 把它们变成同一套接口、
        统一了网格与求解器、并补齐了验证。
\end{itemize}

% ===========================================================================
\chapter{扩展框架}
\label{chap:extend}

\section{新增一种单元}

步骤：

\begin{enumerate}[label=(\arabic*),itemsep=0.2em]
  \item 新建 \code{MosicaFE/spaces/yourspace.py}，继承 \code{BaseFESpace}；
  \item 实现 \code{\_build\_dof\_map()}：建立单元局部自由度到全局自由度的映射，
        并填写 \code{n\_dofs}、\code{dof\_support}、\code{dof\_is\_interior}；
  \item 实现 \code{stiffness\_matrix()}、\code{mass\_matrix()}、\code{load\_vector(f)}；
  \item 实现 \code{interpolate(f)} 与 \code{evaluate(cell, u, points)}；
  \item 在 \code{spaces/factory.py} 的 \code{FESpace.FAMILIES} 中注册；
  \item 在 \code{tests/} 中补上 patch test 与收敛阶测试。
\end{enumerate}

\begin{lstlisting}[language=MosicaFE,caption={示例 12.1：显式基函数空间的骨架},label={lst:newspace}]
from MosicaFE.spaces.base import BaseFESpace

class MySpace(BaseFESpace):
    family = "myspace"
    has_explicit_basis = True

    def _build_dof_map(self):
        # fill self.dof_map, self.n_dofs, self.dof_support, self.dof_is_interior
        ...

    def reference_nodes(self, cell_idx):
        # reference coordinates of local DOFs
        ...

    def eval_basis(self, xi):
        ...

    def eval_basis_grad(self, xi):
        ...

    def stiffness_matrix(self):
        return self._assemble(self._local_stiffness, symmetric=True)

    def mass_matrix(self):
        return self._assemble(self._local_mass, symmetric=True)

    def load_vector(self, f):
        return self._assemble_vector(lambda c: self._local_load(c, f))

    def interpolate(self, f):
        return self._evaluate_function(f, self.dof_coords)
\end{lstlisting}

\begin{important}
如果新单元的基函数是显式的（例如更高次的 Lagrange、DG、
或张量积的其它族），可以参考 \code{spaces/lagrange.py}：
它已经实现了通用的"参考单元求值 + 逐单元变换"组装循环，
通常只需要提供 \code{eval\_basis} 与自由度布局。
\end{important}

\section{新增一门方程}

\begin{lstlisting}[language=MosicaFE,caption={示例 12.2：新 PDE 模板},label={lst:newpde}]
from MosicaFE.physics.base import BasePhysics
from MosicaFE.physics.registry import register_problem

@register_problem("my_pde")
class MyProblem(BasePhysics):
    """Weak form: a(u, v) = l(v)."""

    def __init__(self, space, f=None, g=None, **kwargs):
        super().__init__(space, f=f, g=g, **kwargs)
        # store problem specific data here

    def assemble(self):
        K = self.space.stiffness_matrix()
        M = self.space.mass_matrix()
        A = K + 3.0 * M                      # example combination
        b = self.space.load_vector(self.f)
        return self.apply_boundary_conditions(A, b)
\end{lstlisting}

注册之后，\code{available\_problems()} 会包含 \code{"my\_pde"}，
并且 \code{create\_problem("my\_pde", V, ...)} 与
\code{run("my\_pde", mesh, element="VEM", ...)} 立刻可用。

\section{新增一种混合格式（以 RT0$\times$P0 为模板）}

混合格式（鞍点问题）与上面的"椭圆算子"套路不同，但分工是一样的：
\textbf{空间算出单元矩阵块，物理层只做拼装}。
RT0$\times$P0 在库内就是按这个模板实现的，可以照抄：

\begin{enumerate}[label=(\arabic*),itemsep=0.2em]
  \item 如果需要新的面自由度布局，先在 \code{core/facets.py} 拿到面拓扑
        （全局面编号、\code{cell\_facets}、\code{signs}）；
  \item 在空间里实现"块"接口：本例是
        \code{flux\_mass\_matrix()}（$(\Phi,\Phi)$）与
        \code{divergence\_matrix()}（$(P,\Phi)$），
        外加 \code{pressure\_load\_vector(f)}；
  \item 在物理层用 \code{bmat}（来自 \code{MosicaFE.\_compat}）把块拼成鞍点矩阵：
        \code{bmat([[M, -B.T], [B, None]])}；
  \item 求解方式：\code{"direct"}、\code{"schur"}（静力凝聚）
        或 \code{Solver.minres}（对称不定）；
  \item 误差分析放进 \code{postprocess/error.py}，
        参照 \code{mixed\_error\_norms}（压力 + 通量各一套范数）。
\end{enumerate}

\begin{important}
\code{physics} 层只有 6 行拼装代码（见
\code{physics/mixed\_poisson.py} 的 \code{assemble}），
而单元几何与数值积分全部留在 \code{spaces/rt0.py}。
因此换一套混合元（例如 BDM1、或 Stokes 的 $Q_2\times Q_1$ 配对）
时，只需要提供新的块，物理层与后处理可以原样复用。
\end{important}

\section{新增一种网格}

\begin{enumerate}[label=(\arabic*),itemsep=0.2em]
  \item 在 \code{core/mesh.py} 中新增一个类方法，
        返回 \code{Mesh(nodes, elements, element\_type, name=...)}；
  \item 2D 单元会自动按带符号面积纠正为逆时针；
        3D 四面体会按体积纠正取向；
  \item 如果 3D 单元不是四面体 / 六面体，
        请通过 \code{faces=...} 提供每个单元的面表（局部顶点编号的列表）。
\end{enumerate}

\begin{lstlisting}[language=MosicaFE,caption={示例 12.3：自定义单元形状}]
import numpy as np
from MosicaFE import Mesh

# a two-element hexagonal domain (arbitrary polygons are fine)
nodes = np.array([[0.0, 0.0], [1.0, 0.0], [1.5, 0.9],
                  [0.5, 1.5], [-0.5, 0.9], [0.5, 0.6]])
elements = [[0, 1, 2, 5], [0, 5, 3, 4],
            [1, 2, 3, 5], [0, 3, 4, 5]]

mesh = Mesh.from_arrays(nodes, elements, boundary_nodes=[0, 1, 2, 3, 4])
print(mesh.summary())
\end{lstlisting}

\section{开发流程与测试要求}

建议的流程：

\begin{enumerate}[label=(\arabic*),itemsep=0.2em]
  \item 先写\textbf{失败}的测试（patch test、收敛阶、对称性等）；
  \item 再实现功能；
  \item 跑 \code{tests/run\_tests.py} 与 \code{tests/run\_checks.py}；
  \item 更新 \code{progress/} 下的进度与验证记录。
\end{enumerate}

每个新特性至少要覆盖：

\begin{itemize}[itemsep=0.2em]
  \item 单元级：基函数的 Kronecker 性质、单位分解、单元矩阵对称性/正定性；
  \item 空间级：全局矩阵的对称性、常数核、$M$ 行和与 $b(f\equiv1)$ 的一致性；
  \item 问题级：patch test（线性精确解必须精确）、收敛阶；
  \item 端到端：\code{solve\_poisson} / \code{run} 一行接口可用。
\end{itemize}

% ===========================================================================
\chapter{常见问题与排查}
\label{chap:faq}

\section{常见问题}

\subsection{导入错误}

\textbf{现象}：\code{ModuleNotFoundError: No module named 'MosicaFE'}。

\textbf{原因}：\mosica{} 在 \code{D:\textbackslash mosica\textbackslash MosicaFE}，
需要把\textbf{父目录} \code{D:\textbackslash mosica} 加入模块路径。

\begin{lstlisting}[style=shell]
# option 1: install in editable mode
cd D:\mosica\MosicaFE && python -m pip install -e .

# option 2: extend PYTHONPATH
$env:PYTHONPATH = "D:\mosica"
\end{lstlisting}

\subsection{收敛阶不对}

排查顺序：

\begin{enumerate}[label=(\arabic*),itemsep=0.15em]
  \item 先跑 patch test（第 \ref{chap:physics} 章的式 \eqref{eq:patch}）。
        若不通过，问题一定在组装或边界处理；
  \item 检查求积阶是否足够（张量积单元尤其容易欠积分）；
  \item 检查 \(h\) 的定义是否与网格加密一致；
  \item 检查精确解是否足够光滑（制造解最好用三角函数或多项式）；
  \item 对 VEM，检查稳定化选项是否为 \code{"trace"}。
\end{enumerate}

\subsection{迭代法不收敛}

\begin{itemize}[itemsep=0.2em]
  \item 先用直接法验证系统本身可解；
  \item CG 要求矩阵对称正定：鞍点问题必须换 MINRES 或直接法；
  \item 增大 \code{maxiter}、放宽 \code{tol}；
  \item 条件数很大时考虑缩放或预处理。
\end{itemize}

\subsection{内存不足}

\begin{itemize}[itemsep=0.2em]
  \item 确认在使用稀疏矩阵（安装了 \code{scipy}）；
  \item 三维问题网格数增长很快：$n^3$ 个节点对应约 $6n^3$ 个四面体；
  \item 改用迭代法可避免直接法的填充（fill-in）。
\end{itemize}

\section{常见问答}

\paragraph{问：怎样选择多项式次数？}
答：绝大多数情形 P2 / Q2 的精度/开销比最好；先用 P1 做快速验证，
解很光滑时再上二次元。虚拟元目前只有 $k=1$。

\paragraph{问：什么时候用迭代法？}
答：自由度少于约 $2\times10^4$ 时直接法最省心；
更大规模改用 CG（对称正定）或 GMRES（一般）。

\paragraph{问：VEM 与 FEM 到底差在哪？}
答：在三角形网格上，$k=1$ 虚拟元空间恰好等于 P1 有限元空间，
两者数值结果完全相同（本库实测两者误差逐位一致）。
VEM 的优势出现在\textbf{一般多边形网格}上：它允许每个单元的面数不同、
允许非结构的多边形剖分，而经典 Lagrange 有限元需要单元形状统一。

\paragraph{问：为什么 RT0 的解返回两个数组？}
答：混合格式同时求解通量 $q_h$ 与压力 $p_h$。
\code{Q[c,i]} 是单元 $c$ 第 $i$ 个面（按"不含顶点 $i$"编号）上的
外法向通量，\code{P[c]} 是该单元的常数压力。

\paragraph{问：RT0 需要安装或接入外部程序吗？}
答：不需要。RT0$\times$P0 已经是 \mosica{} 自己的方法：显式基函数
$\varphi_i=(x-a_i)/(d|K|)$ 与精确质量矩阵在 \code{spaces/rt0.py}，
面拓扑与定向在 \code{core/facets.py}，鞍点组装在
\code{physics/mixed\_poisson.py}，误差范数在 \code{postprocess/error.py}。
库的依赖仍然只有 numpy（scipy / matplotlib 可选），
\code{import MosicaFE} 之后 \code{FESpace.rt0(mesh)} 直接可用。

\paragraph{问：能求解时间相关问题吗？}
答：当前版本只处理稳态椭圆问题。时间相关问题建议用"线方法"：
空间离散得到 $M\dot u+Au=b$，时间方向自行用显式/隐式格式推进。
\code{mass\_matrix()} 已经提供。

\paragraph{问：能求解非线性问题吗？}
答：框架面向线性问题。非线性问题建议在外层写 Newton 迭代，
每一步用 \mosica{} 组装雅可比并求解。

\paragraph{问：怎样把结果导出到 ParaView？}
答：用 \code{space.nodal\_values(u)} 取得节点值，
配合 \code{mesh.nodes} 与 \code{mesh.elements} 自行写
\code{.vtu} 或 \code{.vtk}；也可以在 \code{postprocess} 中新增导出函数。

\appendix

% ===========================================================================
\chapter{API 速查}
\label{app:api}

\section{顶层模块 \module{}}

\begingroup\footnotesize
\begin{longtable}{@{}L{5.0cm}L{9.4cm}@{}}
\toprule
名字 & 说明 \\
\midrule
\endfirsthead
\toprule
名字 & 说明 \\
\midrule
\endhead
\code{Mesh} & 网格对象与工厂方法 \\
\code{FESpace} & 空间工厂：\code{lagrange} / \code{vem} / \code{rt0} / \code{create} \\
\code{FacetTopology}, \code{build\_facet\_topology} & 面拓扑与定向（$H(\mathrm{div})$ 空间） \\
\code{QuadratureRule} & 参考单元求积规则 \\
\code{cell\_quadrature(...)} / \code{simplex\_quadrature(...)} & 物理单元 / 任意单纯形上的求积点与权重 \\
\code{BaseFESpace} & 空间基类 \\
\code{LagrangeSpace} / \code{VEMSpace} / \code{RT0Space} & 三种空间 \\
\code{rt0\_flux\_basis}, \code{rt0\_mass\_matrices} & RT0 显式基函数与精确质量矩阵 \\
\code{BasePhysics} & 物理问题基类 \\
\code{PoissonProblem} & $-\nabla\cdot(\kappa\nabla u)=f$ \\
\code{MixedPoissonProblem} & RT0$\times$P0 混合形式 \\
\code{ReactionDiffusionProblem} & $-\Delta u+cu=f$ \\
\code{register\_problem}, \code{create\_problem}, \code{available\_problems} & PDE 注册表 \\
\code{Solver} & 线性求解器（direct / cg / gmres / bicgstab / minres / auto） \\
\code{compute\_error} / \code{compute\_errors} & 误差范数 \\
\code{mixed\_error\_norms} & RT0$\times$P0 的压力 / 通量误差范数 \\
\code{convergence\_study} & 收敛性研究 \\
\code{estimate\_convergence\_rate} & 收敛阶估计 \\
\code{plot\_solution} / \code{plot\_mesh} / \code{plot\_error} & 可视化 \\
\code{plot\_convergence} / \code{plot\_pressure\_2d} / \code{plot\_flux\_2d} & 可视化 \\
\code{solve\_poisson} & 一行式 Poisson 求解 \\
\code{run} & 按名字求解任意已注册 PDE \\
\bottomrule
\end{longtable}
\endgroup

\section{\code{Mesh}}

\begingroup\small
\begin{longtable}{@{}L{5.0cm}L{9.4cm}@{}}
\toprule
方法 / 属性 & 说明 \\
\midrule
\endfirsthead
\toprule
方法 / 属性 & 说明 \\
\midrule
\endhead
\code{Mesh.interval(a,b,nx)} & 一维区间网格 \\
\code{Mesh.rectangle(...)} & 二维矩形网格（\code{triangle}/\code{quad}） \\
\code{Mesh.box(...)} & 三维长方体网格（\code{tet}/\code{hex}） \\
\code{Mesh.pentagon(nx,ny)} & 五边形网格 \\
\code{Mesh.voronoi\_polygon(...)} & Voronoi 多边形网格 \\
\code{Mesh.from\_arrays(...)} & 自定义网格 \\
\code{mesh.refine\_uniform()} & 均匀加密 \\
\code{mesh.get\_mesh\_size()} & 最大单元直径 $h$ \\
\code{mesh.get\_cell\_vertices(i)} & 单元顶点坐标 \\
\code{mesh.get\_cell\_measure(i)} & 单元面积 / 体积 \\
\code{mesh.get\_cell\_diameter(i)} & 单元直径 \\
\code{mesh.get\_cell\_centroid(i)} & 单元形心 \\
\code{mesh.boundary\_nodes} & 边界顶点编号 \\
\code{mesh.boundary\_facets} / \code{mesh.facets} & facet 拓扑 \\
\code{mesh.facet\_topology} & \code{FacetTopology}：全局面编号、共享面、外法向符号 \\
\code{mesh.summary()} & 一行摘要 \\
\bottomrule
\end{longtable}
\endgroup

\section{\code{FESpace} 与空间方法}

\begingroup\small
\begin{longtable}{@{}L{5.0cm}L{9.4cm}@{}}
\toprule
方法 & 说明 \\
\midrule
\endfirsthead
\toprule
方法 & 说明 \\
\midrule
\endhead
\code{FESpace.lagrange(mesh, degree)} & 连续 Lagrange（P1/P2/Q1/Q2） \\
\code{FESpace.vem(mesh, degree=1)} & 虚拟元（$k=1$） \\
\code{FESpace.rt0(mesh)} & RT0$\times$P0 混合空间 \\
\code{FESpace.create(mesh, family, degree)} & 按名字创建 \\
\code{FESpace.available()} & 可用空间族 \\
\code{V.n\_dofs} & 自由度个数 \\
\code{V.get\_cell\_dofs(i)} & 单元局部 → 全局自由度 \\
\code{V.get\_dof\_coordinates(k)} & 自由度坐标 \\
\code{V.boundary\_dofs()} & Dirichlet 自由度 \\
\code{V.eval\_basis(xi)} / \code{V.eval\_basis\_grad(xi)} & 参考单元基函数 \\
\code{V.stiffness\_matrix()} & 刚度矩阵 \\
\code{V.mass\_matrix()} & 质量矩阵 \\
\code{V.load\_vector(f)} & 载荷向量 \\
\code{V.interpolate(f)} & 插值 \\
\code{V.evaluate(c, u, pts)}, \code{V.evaluate\_gradient(...)} & 单元内重构求值 \\
\code{V.nodal\_values(u)} & 节点值（可视化用） \\
\midrule
\code{V.flux\_mass\_matrix()} & RT0：$(\Phi,\Phi)$ 块（精确闭式质量矩阵） \\
\code{V.divergence\_matrix()} & RT0：$(P,\Phi)$ 块（散度，恒为 $\pm1$） \\
\code{V.pressure\_load\_vector(f)} & RT0：压力方程右端项 $\int_K f$ \\
\code{V.interpolate\_pressure(f)} / \code{V.interpolate\_flux(q)} & RT0：两个场各自的插值 \\
\code{V.evaluate\_flux(c, Q, pts)} & RT0：单元内重构 $q_h$ \\
\code{V.locate\_cell(x)} & RT0：定位点所在单元 \\
\code{V.n\_flux} / \code{V.n\_pressure} / \code{V.topology} & RT0：自由度规模与面拓扑 \\
\bottomrule
\end{longtable}
\endgroup

\section{物理问题}

\begingroup\small
\begin{longtable}{@{}L{5.0cm}L{9.4cm}@{}}
\toprule
方法 & 说明 \\
\midrule
\endfirsthead
\toprule
方法 & 说明 \\
\midrule
\endhead
\code{PoissonProblem(V, f, g, kappa)} & 构造 \\
\code{problem.assemble()} & 返回 $(A,b)$（含边界条件） \\
\code{problem.solve(method)} & 组装并求解 \\
\code{problem.compute\_error(u, exact, norm)} & 单个误差范数 \\
\code{problem.compute\_errors(u, exact, grad)} & 全部误差范数 \\
\code{MixedPoissonProblem(...)} & 返回 $(Q,P)$ \\
\code{mixed.error\_norms(P, Q, p\_exact, q\_exact)} & 压力/通量误差 \\
\code{mixed.check\_conservation()} & 单元质量守恒误差 \\
\code{mixed.flux\_at\_points(points)} & 给定点上的通量重构（自动定位单元） \\
\code{mixed.cell\_flux(cell)} & 单元形心处的通量向量（画箭头用） \\
\code{mixed.solve(method)} & 鞍点系统的三种求解方式：\code{direct} / \code{schur} / \code{minres} \\
\bottomrule
\end{longtable}
\endgroup

\section{求解器与后处理}

\begingroup\small
\begin{longtable}{@{}L{5.0cm}L{9.4cm}@{}}
\toprule
方法 & 说明 \\
\midrule
\endfirsthead
\toprule
方法 & 说明 \\
\midrule
\endhead
\code{Solver.direct(A,b)} & 直接法 \\
\code{Solver.conjugate\_gradient(...)} & CG \\
\code{Solver.gmres(A,b,tol)} & GMRES \\
\code{Solver.bicgstab(A,b,tol)} & BiCGSTAB \\
\code{Solver.minres(A,b,tol)} & MINRES \\
\code{Solver.solve(A,b,method)} & 统一入口 \\
\code{Solver.apply\_dirichlet(...)} & 强加 Dirichlet \\
\code{Solver.recommend(n, symmetric)} & 推荐求解器 \\
\code{convergence\_study(...)} & 收敛性研究 \\
\code{estimate\_convergence\_rate(h, e)} & 收敛阶 \\
\code{plot\_solution(u, mesh, V, ...)} & 画解 \\
\code{plot\_mesh(mesh, ...)} & 画网格 \\
\code{plot\_error(u, mesh, space, exact)} & 画误差 \\
\code{plot\_convergence(h, errors)} & 画收敛曲线 \\
\code{plot\_pressure\_2d(mesh, P)} & 画 RT0 压力 \\
\code{plot\_flux\_2d(space=V, Q=Q)} & 画通量（用 $q_h$ 在单元形心处重构） \\
\code{compare\_solutions(entries)} & 并排比较 \\
\bottomrule
\end{longtable}
\endgroup

% ===========================================================================
\chapter{数学符号}
\label{app:notation}

\section{一般记号}

\begin{longtable}{@{}ll@{}}
\toprule
符号 & 含义 \\
\midrule
\endfirsthead
\toprule
符号 & 含义 \\
\midrule
\endhead
$\Omega$ & 求解区域 \\
$\partial\Omega$ & 区域边界 \\
$d$ & 空间维数（1、2、3） \\
$h$ & 网格尺寸（最大单元直径） \\
$K$ & 网格单元 \\
$\hat K$ & 参考单元 \\
$n$ & 自由度个数 \\
$\mathcal{D}$ & Dirichlet 自由度集合 \\
\bottomrule
\end{longtable}

\section{函数空间}

\begin{longtable}{@{}ll@{}}
\toprule
符号 & 含义 \\
\midrule
\endfirsthead
\toprule
符号 & 含义 \\
\midrule
\endhead
$V$ & 无穷维试探/检验空间 \\
$V_h\subset V$ & 有限维离散空间（有限元或虚拟元） \\
$H^1(\Omega)$ & 一阶 Sobolev 空间 \\
$H(\mathrm{div};\Omega)$ & 散度平方可积的向量场空间 \\
$L^2(\Omega)$ & 平方可积函数空间 \\
$\mathcal{P}_k$ & 次数不超过 $k$ 的多项式空间 \\
$Q_k$ & 每个方向次数不超过 $k$ 的张量积多项式空间 \\
\bottomrule
\end{longtable}

\section{算子与范数}

\begin{longtable}{@{}ll@{}}
\toprule
符号 & 含义 \\
\midrule
\endfirsthead
\toprule
符号 & 含义 \\
\midrule
\endhead
$\nabla u$ & 梯度 \\
$\nabla\cdot q$ & 散度 \\
$\Delta u$ & Laplace 算子 \\
$\varphi_i$ / $\hat\varphi_i$ & 物理 / 参考单元上的基函数 \\
$\Pi^{\nabla}$ & VEM 的梯度投影算子 \\
$\Pi_{\rm dof}$ & 投影算子在自由度上的矩阵表示 \\
$A_{\rm proj}$, $S_{\rm stab}$ & VEM 单元矩阵的一致性项与稳定化项 \\
$J_K$ & 参考单元到物理单元的雅可比矩阵 \\
$\|e\|_{L^2}$, $|e|_{H^1}$, $\|e\|_{L^\infty}$ & 误差范数 \\
$q$, $p$ & RT0 混合格式中的通量与压力 \\
$A$, $b$ & 离散线性系统的矩阵与右端项 \\
\bottomrule
\end{longtable}

% ===========================================================================
\chapter{参考文献}
\label{app:bib}

\begin{thebibliography}{9}
\bibitem{brenner} S. C. Brenner and L. R. Scott,
  \emph{The Mathematical Theory of Finite Element Methods}, 3rd ed.,
  Springer, 2008.
\bibitem{ern} A. Ern and J.-L. Guermond,
  \emph{Theory and Practice of Finite Elements}, Springer, 2004.
\bibitem{logg} A. Logg, K.-A. Mardal, and G. Wells (eds.),
  \emph{Automated Solution of Differential Equations by the Finite Element Method},
  Springer, 2012.
\bibitem{hitchhiker} L. Beir\~ao da Veiga, F. Brezzi, A. Cangiani,
  G. Manzini, L. D. Marini, and A. Russo,
  ``Basic principles of virtual element methods'',
  \emph{Mathematical Models and Methods in Applied Sciences}, 23(1):199--214, 2013.
\bibitem{vem2014} L. Beir\~ao da Veiga, F. Brezzi, L. D. Marini, and A. Russo,
  ``The Hitchhiker's Guide to the Virtual Element Method'',
  \emph{Mathematical Models and Methods in Applied Sciences}, 24(8):1541--1573, 2014.
\bibitem{vemstab} L. Beir\~ao da Veiga, F. Brezzi, L. D. Marini, and A. Russo,
  ``Virtual element methods for general second order elliptic problems on
  polygonal meshes'', \emph{SIAM Journal on Numerical Analysis}, 2013.
\bibitem{brezzi} F. Brezzi and M. Fortin,
  \emph{Mixed and Hybrid Finite Element Methods}, Springer, 1991.
\bibitem{rt} P.-A. Raviart and J.-M. Thomas,
  ``A mixed finite element method for 2nd order elliptic problems'',
  in \emph{Mathematical Aspects of Finite Element Methods}, Springer, 1977.
\bibitem{saad} Y. Saad, \emph{Iterative Methods for Sparse Linear Systems},
  2nd ed., SIAM, 2003.
\end{thebibliography}

\end{document}
